\chapter{Trigonometrie}
\setcounter{figure}{1}
\setcounter{subfigure}{1}

Trigonometrie (uitgespreek tri-go-no-me-trie) handel oor die verwantskappe tusen die hoeke en die sye van 'n reghoekige driehoek. Ons sal leer van trigonometriese funksies wat die basis vorm van trigonometrie.\par 

\section{Trigonometrie is bruikbaar}
\nopagebreak
Daar is vele toepassings van trigonometrie. Van spesifieke belang is die tegniek van driehoeksmeting, wat gebruik word in sterrekunde om die afstand na nabye sterre te meet, in aardrykskunde (geografie) om die afstand tussen landmerke te meet en in satellietnavigasie. GPS (globale posisioneringstelsels) sou nie moontlik gewees het sonder trigonometrie nie. Ander gebiede wat trigonometrie gebruik sluit in akoestiek, optika, analise van finansiele markte, elektronika, waarskynlikheidsteorie, statistiek, biologie, mediese grafiese werk (CAT skandering en ultraklank), chemie, kriptologie, meteorologie, oseanografie, landmeetkunde, argitektuur, fonetika, ingenieurswese, rekenaargrafika, en in die ontwikkeling van rekenaarspeletjies.\par 

% I moved this to here in die english version too. Methinks it belongs here
% horizontal\label{m39405*id78218}% \label{m39405*secfhsst!!!underscore!!!id95}
%             \subsection{  Bespreking : Gebruike van Trigonometrie }
%             \nopagebreak
%             
%       \label{m39405*id78126}Kies een van die velde waar trigonometrie gebruik word uit die lys hierbo en skryf 'n verslag van 1 bladsy wat beskryf \textsl{hoe} trigonometrie in die gekose veld gebruik word. \par 

\section{Gelykvormigheid van driehoeke}

As $\triangle ABC$ gelykvormig is aan $ \triangle DEF$, dan word dit geskryf as:\par 

\begin{equation*}
\triangle ABC||| \triangle DEF
\end{equation*}

\setcounter{subfigure}{0}
\begin{figure}[H] 
\begin{center}
\begin{pspicture}(-1,-0.6)(6.6,3)
%\psgrid[gridcolor=gray]
\rput(1.2,0){\pstTriangle(0,0){A}(2; 25){B}(3; 125){C}}
\rput(5.2,0){\pstTriangle[unit=0.5](0,0){D}(2; 25){E}(3; 125){F}}
\end{pspicture}
\end{center}   
\end{figure}   
\par 
Nou is dit moontlik om die verhoudings af te lei tussen ooreenkomstige sye van die twee gelykvormige driehoeke:\par 


\begin{equation*}
\begin{array}{ccl}\hfill \dfrac{AB}{BC}& =& \dfrac{DE}{EF}\hfill \\[5pt]
 \hfill \dfrac{AB}{AC}& =& \dfrac{DE}{DF}\hfill \\ [5pt]
\hfill \dfrac{AC}{BC}& =& \dfrac{DF}{EF}\hfill \\[5pt]
 \hfill \dfrac{AB}{DE}& =& \dfrac{BC}{EF}=\dfrac{AC}{DF}\hfill \end{array}
\end{equation*}
Nog 'n belangrike feit oor die gelykvormige driehoeke $ABC$ and $DEF$ is dat hoek $A$ gelyk is aan hoek $D$, hoek $B$ gelyk is aan hoek $E$ en hoek $C$ is gelyk aan hoek $F$.\par 


\begin{equation*}
\begin{array}{ccc}\hfill \hat A& =& \hat D\hfill \\ \hfill \hat B& =& \hat E\hfill \\ \hfill \hat C& =& \hat F\hfill \end{array}
\end{equation*}

\begin{Investigation}{Verhoudings van gelykvormige driehoeke}

Teken drie gelykvormige driehoeke van verskillende grooottes, met elk $\hat{A}={30}^{\circ }$; $\hat{B}={90}^{\circ }$ en $\hat{C}={60}^{\circ }$. Meet hoeke en lengtes akkuraat en skryf die metings in die tabel (rond die metings af na 1 desimaal).\par 

\setcounter{subfigure}{0}
\scalebox{1} % Change this value to rescale die drawing.
{
\begin{pspicture}(0,-2.3817186)(8.567187,2.3817186)
\psline[linewidth=0.04cm](0.301875,-1.9001563)(1.942836,-1.9001563)
\psline[linewidth=0.04cm](1.942836,-1.9001563)(1.942836,0.84772635)
\psline[linewidth=0.04cm](0.301875,-1.9001563)(1.942836,0.8604187)
\psline[linewidth=0.04cm](2.641875,-1.9201562)(4.641875,-1.9201562)
\psline[linewidth=0.04cm](4.641875,-1.9201562)(4.641875,1.5438437)
\psline[linewidth=0.04cm](2.641875,-1.9201562)(4.641875,1.5598438)
\psline[linewidth=0.04cm](5.521875,-1.9201562)(7.801521,-1.9201562)
\psline[linewidth=0.04cm](7.801521,-1.9201562)(7.801521,1.9617095)
\psline[linewidth=0.04cm](5.521875,-1.9201562)(7.801521,1.9796396)
\usefont{T1}{ptm}{m}{n}
\rput(0.26546875,-2.1101563){$C$}
\usefont{T1}{ptm}{m}{n}
\rput(2.8054688,-2.1901562){$C'$}
\usefont{T1}{ptm}{m}{n}
\rput(5.6254687,-2.1901562){$C''$}
\usefont{T1}{ptm}{m}{n}
\rput(2.0454688,-2.1301563){$B$}
\usefont{T1}{ptm}{m}{n}
\rput(4.7654686,-2.1901562){$B'$}
\usefont{T1}{ptm}{m}{n}
\rput(8.105469,-2.2101562){$B''$}
\usefont{T1}{ptm}{m}{n}
\rput(2.1254687,1.0898438){$A$}
\usefont{T1}{ptm}{m}{n}
\rput(4.865469,1.8098438){$A'$}
\usefont{T1}{ptm}{m}{n}
\rput(8.105469,2.1698437){$A''$}
\psframe[linewidth=0.04,dimen=outer](1.961875,-1.5601562)(1.601875,-1.9201562)
\psframe[linewidth=0.04,dimen=outer](4.661875,-1.5801562)(4.301875,-1.9401562)
\psframe[linewidth=0.04,dimen=outer](7.801875,-1.5801562)(7.441875,-1.9401562)
\usefont{T1}{ptm}{m}{n}
\rput(1.64,-0.25){$30^{\circ}$}
% \pscircle[linewidth=0.0060,dimen=outer](1.811875,0.02984375){0.03}
\usefont{T1}{ptm}{m}{n}
\rput(0.9,-1.6){$60^{\circ}$}
% \pscircle[linewidth=0.0060,dimen=outer](0.951875,-1.6101563){0.03}
\usefont{T1}{ptm}{m}{n}
\rput(3.3,-1.6){$60^{\circ}$}
% \pscircle[linewidth=0.0060,dimen=outer](3.291875,-1.6301563){0.03}
\usefont{T1}{ptm}{m}{n}
\rput(6.2,-1.6){$60^{\circ}$}
% \pscircle[linewidth=0.0060,dimen=outer](6.131875,-1.6301563){0.03}
\usefont{T1}{ptm}{m}{n}
\rput(4.3173437,0.4){$30^{\circ}$}
% \pscircle[linewidth=0.0060,dimen=outer](4.531875,0.64984375){0.03}
\usefont{T1}{ptm}{m}{n}
\rput(7.4973435,0.9){$30^{\circ}$}
% \pscircle[linewidth=0.0060,dimen=outer](7.711875,1.1498437){0.03}
\end{pspicture} 
}      
\par 
% \textbf{m39405*id78604}\par
\begin{table}[H]
% \begin{table}[H]
% \\ '' '0'
\begin{center}

\noindent
\setlength{\extrarowheight}{2pt}

\begin{tabular}{|m{2.5cm}|m{2.5cm}|m{2.5cm}|}\hline
% My position: 0
% my spanname: 
% my ct of spanspec: 0
% my column-count: 3
\multicolumn{3}{|c|}{Deel lengtes van sye (verhoudings) }
\\ \hline
%--------------------------------------------------------------------
\LARGE$\frac{AB}{BC}=$
&
\LARGE$\frac{AB}{AC}=$
&
\LARGE$\frac{CB}{AC}=$
% make-rowspan-placeholders
\\ \hline
%--------------------------------------------------------------------
\LARGE$\frac{A'B'}{B'C'}=$
&
\LARGE$\frac{A'B'}{A'C'}=$
&
\LARGE$\frac{C'B'}{A'C'}=$
% make-rowspan-placeholders
\\ \hline
%--------------------------------------------------------------------
\LARGE$\frac{A''B''}{B''C''}=$
&
\LARGE$\frac{A''B''}{A''C''}=$
&
\LARGE$\frac{C''B''}{A''C''}=$
% make-rowspan-placeholders
\\ \hline
%--------------------------------------------------------------------
\end{tabular}
\end{center}
% \begin{center}{\small\bfseries Table 14.1}\end{center}
% \begin{caption}{\small\bfseries Table 14.1}\end{caption}
\end{table}
\par
Watter waarnemings kan jy maak oor die verhoudings van die sye?\\
\\
Hierdie gelyke verhoudings word gebruik om die trigonometriese verhoudings te definieer.
\end{Investigation}


    

\section{Definisie van die trigonometriese verhoudings}
Beskou 'n reghoekige driehoek $ABC$.\par 

\setcounter{subfigure}{0}
\begin{figure}[H] % horizontal\label{m39408*id79667}
\begin{center}
\begin{pspicture}(-1,-0.6)(4,3)
%\psgrid[gridcolor=gray]
\pstTriangle(0,0){A}(2.4;90){B}(3.35;0){C}
\pstRightAngle{B}{A}{C}
\pstMarkAngle{B}{C}{A}{$\theta$}
\pcline[linestyle=none](A)(B)
\aput{:U}{opposite}
\pcline[linestyle=none](B)(C)
\aput{:U}{hypotenuse}
\pcline[linestyle=none](A)(C)
\bput{:U}{adjacent}
\end{pspicture}
\end{center}
\end{figure}       
\par 

\Note{In algebra gebruik ons gereeld die letter $x$ om die onbekende veranderlike aan te dui (al kan ons enige ander letter soos $a$, $b$, $k$, ens. ook gebruik. In trigonometrie gebruik ons gereeld die Griekse letter $\theta $ vir 'n onbekende hoek (ons gebruik ook $\alpha $, $\beta $, $\gamma $ ens).}


In die reghoekige driehoek vrewys ons na die lengtes van die sye na gelang van hulle plasing in verhouding met die hoek $\theta $. Die sy oorkant die regte hoek is die skuinssy, die sy oorkant $\theta $ is die teenoorstaande sy en die sy naas $\theta $ is die aanliggende sy.
Let op dat die keuse van die nie-$90^{\circ}$ binnehoek arbitrer is. Jy kan enige van die binnehoeke kies en dan die aanliggende en teenoorstaande sy definieer. Die skuinssy bly egter die skuinssy ongeag die binnehoek van belang (omdat dit ALTYD oorkant die regte hoek is en ALTYD die langste sy).\par 

Ons definieer die trigonometriese verhoudings sinus, cosinus en tangens soos volg:
\par 

\begin{equation*}
\begin{array}{ccc}\hfill sin~\theta & =& \frac{\mbox{opposite}}{\mbox{hypotenuse}}\hfill \\
\\
 \hfill cos~\theta & =& \frac{\mbox{adjacent}}{\mbox{hypotenuse}}\hfill \\
\\
 \hfill tan~\theta & =& \frac{\mbox{opposite}}{\mbox{adjacent}}\hfill 
\end{array}
\end{equation*}

Hierdie verhoudings, ook bekend as die trigonometriese identiteite, bepaal die verhoudings van die lengtes van 'n driehoek se sye tot sy binnehoeke.\par 

Hierdie drie verhoudings vorm die basis van trigonometrie. \par
% Hierdie volgende rympie-voorbeeld is nie vertaalbaar na Afrikaans nie, ek stel voor julle haal dit uit die Engels.
You may also hear people saying ``Soh Cah Toa''. This is a mnemonic technique for remembering die trigonometric ratios:\par 

% \textbf{m39408*id79953}\par
\begin{table}[H]
% \begin{table}[H]
% \\ '' '0'
\begin{center}
\label{m39408*id79953}
\noindent

\begin{tabular}{|l|}\hline

$\mathbf{s}in=\frac{\mbox{\textbf{o}pposite}}{\mbox{\textbf{h}ypotenuse}} $
% make-rowspan-placeholders

\\  \hline
%--------------------------------------------------------------------

$\mathbf{c}os=\frac{\mbox{\textbf{a}djacent}}{\mbox{\textbf{h}ypotenuse}} $
% make-rowspan-placeholders

\\ \hline
%--------------------------------------------------------------------

$\mathbf{t}an=\frac{\mbox{\textbf{o}pposite}}{\mbox{\textbf{a}djacent}} $
% make-rowspan-placeholders

\\   \hline
%--------------------------------------------------------------------
\end{tabular}
\end{center}
% \begin{center}{\small\bfseries Table 14.2}\end{center}
% \begin{caption}{\small\bfseries Table 14.2}\end{caption}
\end{table}
\par

\Tip{Die definisies van teenoorstaande, aanliggend en skuinssy is net van toepassing wanneer mens met 'n reghoekige driehoek werk! Bevestig telkens dat jou driehoek 'n regte hoek bevat alvorens jy die verhoudings gebruik, anders kry jy die verkeerde antwoord.}

\begin{activity}{Sakrekenaarwerk}

Dit is baie belangrik dat jy vertroud is met jou sakrekenaar se trigonometrie-knoppies en hoe hulle werk. \\
Gebruik die volgende voorbeelde om te verseker dat jy jou sakrekenaar korrek gebruik:
\begin{enumerate}[noitemsep, label=\textbf{\arabic*}. ] 
 \item $cos~ 30^{\circ} = 0,866$
\item $sin~90^{\circ} = 1$
\item $tan~60^{\circ} = 1,732$
\item $4~sin~45^{\circ}=2,828$
\item $\frac{1}{3}~cos~60^{\circ}=1,5$
\end{enumerate}

\end{activity}

\begin{wex}{Gebruik van jou sakrekenaar}

 {Gebruik jou sakrekenaar om die volgende te bereken (Korrek tot $2$ desimale plekke):
\begin{enumerate}[itemsep=4pt, label=\textbf{\arabic*}. ] 
 \item $cos~48^{\circ}$
\item $2~sin~35^{\circ}$
\item $tan^{2}~81^{\circ}$
\item $3~sin^{2}~72^{\circ}$
\item $\dfrac{1}{4}~cos~27^{\circ} $
\item $\dfrac{5}{6}~tan~34^{\circ}$
\end{enumerate}

}
{

\westep{}
Druk \fbox{COS} \fbox{48} \fbox{\LARGE =} $0,67$

\westep{}
Druk \fbox{2} \fbox{SIN} \fbox{35} \fbox{\LARGE =} $1,15$

\westep{}
Druk \fbox{(} \fbox{TAN} \fbox{81} \fbox{)} \fbox{x$^{2}$}  \fbox{\LARGE =} $39,83$
\\
OF\\
Druk \fbox{TAN} \fbox{81} \fbox{\LARGE =} \fbox{ANS} \fbox{x$^{2}$} \fbox{\LARGE =} $39,83$

\westep{}
Druk \fbox{3} \fbox{(} \fbox{SIN} \fbox{72} \fbox{)} \fbox{x$^{2}$} \fbox{\LARGE =} $2,71$
\\
OF\\
Druk \fbox{SIN} \fbox{72} \fbox{\LARGE =} \fbox{ANS} \fbox{x$^{2}$} \fbox{\LARGE =} \fbox{ANS} \fbox{$\times$} \fbox{3}

\westep{}
Druk \fbox{(} \fbox{1} \fbox{$\div$} \fbox{4} \fbox{)} \fbox{COS} \fbox{27} \fbox{\LARGE =} $0,22$
\\
OF\\
Druk \fbox{COS} \fbox{27} \fbox{\LARGE =} \fbox{ANS} \fbox{$\div$} \fbox{4} \fbox{\LARGE =} $0,22$

\westep{}
Druk \fbox{(} \fbox{5} \fbox{$\div$} \fbox{6} \fbox{)} \fbox{TAN} \fbox{34} \fbox{\LARGE =} $0,56$
\\
OF\\
Druk \fbox{TAN} \fbox{34} \fbox{\LARGE =} \fbox{ANS} \fbox{$\times$} \fbox{5} \fbox{$\div$} \fbox{6} \fbox{\LARGE =} $0,56$

}
\end{wex}
\Note{Die meeste wetenskaplike sakrekenaars is naastenby dieselfde maar hierdie stappe mag verskil afhangende van die sakrekenaar wat jy gebruik. Maak seker jou sakrekenaar is in 'degrees' mode. }

\begin{wex}
{Sakrekenaarwerk}
{As $x=25^{\circ}$ en $y=65^{\circ}$ is, gebruik jousakrekenaar om te bepaal of die volgende stelling waar of vals is:
\begin{equation*}
sin^{2}~x + cos^{2}~(90^{\circ}-y) = 1
\end{equation*}
}
{
\westep{Bereken die linkerkant van die vergelyking}
Druk \fbox{(} \fbox{SIN} \fbox{25} \fbox{)} \fbox{x$^{2}$} \fbox{\LARGE +} \fbox{(} \fbox{COS} \fbox{(} \fbox{90} \fbox{\LARGE -} \fbox{65} \fbox{)} \fbox{)} \fbox{x$^{2}$}  \fbox{\LARGE =} $1$


\westep{Skryf die finale antwoord neer}
LK = RK dus is die stelling waar.



}
\end{wex}

\begin{exercises}{}
{
\begin{enumerate}[itemsep=5pt, label=\textbf{\arabic*}. ]
\item In elk van die volgende driehoeke, dui aan watter van $a$, $b$ en $c$ is die skuinssy, die teenoorstaande sy en die aanliggende sy van die driehoek ten opsigte van $\theta$. 
% \setcounter{subfigure}{0}


\begin{center}
\scalebox{0.85} % Change this value to rescale die drawing.
{
\begin{pspicture}(0,-3.7754679)(20.177345,3.9279697)
\rput (1,2){\textbf{(a)}}
\rput (5.5,2){\textbf{(b)}}
\rput (10.5,2){\textbf{(c)}}
\rput (1.5,-1.5){\textbf{(d)}}
\rput (6,-1.5){\textbf{(e)}}
\rput (9.5,-1.5){\textbf{(f)}}
\psdots[dotsize=0.027999999](2.9173439,-1.5720303)
\psline[linewidth=0.04cm](5.517344,2.5079696)(7.9173436,0.107969664)
\psline[linewidth=0.04cm](7.9173436,0.107969664)(9.0173435,1.2079697)
\psline[linewidth=0.04cm](5.517344,2.5079696)(9.0173435,1.2079697)
\psline[linewidth=0.04cm](7.717344,0.30796966)(7.9173436,0.5079697)
\psline[linewidth=0.04cm](7.9173436,0.5079697)(8.117344,0.30796966)
\psline[linewidth=0.04cm](11.0573435,2.2679696)(11.0573435,0.067969665)
\psline[linewidth=0.04cm](11.0573435,0.067969665)(14.757344,0.067969665)
\psline[linewidth=0.04cm](11.0573435,2.2679696)(14.757344,0.067969665)
\psline[linewidth=0.04cm](1.2373447,-3.3120303)(3.9373438,-0.71203035)
\psline[linewidth=0.04cm](3.9373438,-0.71203035)(3.9373438,-3.3120303)
\psline[linewidth=0.04cm](3.9373438,-3.3120303)(1.2373447,-3.3120303)
\psline[linewidth=0.04cm](7.477345,-1.2120303)(5.7773438,-3.3120303)
\psline[linewidth=0.04cm](7.477345,-1.2120303)(8.9773445,-2.4120302)
\psline[linewidth=0.04cm](8.9773445,-2.4120302)(5.7773438,-3.3120303)
\psline[linewidth=0.04cm](7.257343,-1.5120304)(7.537344,-1.7320304)
\psline[linewidth=0.04cm](7.517344,-1.7120303)(7.717344,-1.4320303)
\psline[linewidth=0.04cm](11.097343,0.38796967)(11.397344,0.38796967)
\psline[linewidth=0.04cm](11.397344,0.38796967)(11.397344,0.08796967)
\psline[linewidth=0.04cm](3.9173427,-3.0320303)(3.6173437,-3.0320303)
\psline[linewidth=0.04cm](3.6173437,-3.0320303)(3.6173437,-3.3320303)
\psline[linewidth=0.04cm](10.0173435,-1.7320304)(11.417343,-3.2320304)
\psline[linewidth=0.04cm](11.417343,-3.2320304)(13.0173435,-1.7320304)
\psline[linewidth=0.04cm](10.0173435,-1.7320304)(13.0173435,-1.7320304)
\psline[linewidth=0.04cm](11.217343,-3.0320303)(11.417343,-2.8320303)
\psline[linewidth=0.04cm](11.417343,-2.8320303)(11.617344,-3.0320303)
\usefont{T1}{ptm}{m}{n}
\rput(8.662344,1.1179696){$\theta$}
\usefont{T1}{ptm}{m}{n}
\rput(13.842344,0.29796967){$\theta$}
\usefont{T1}{ptm}{m}{n}
\rput(3.682343,-1.3020303){$\theta$}
\usefont{T1}{ptm}{m}{n}
\rput(6.362344,-2.9020302){$\theta$}
\usefont{T1}{ptm}{m}{n}
\rput(10.502343,-1.9420303){$\theta$}
\psline[linewidth=0.04cm](2.2973437,2.2879696)(0.09734377,0.18796967)
\psline[linewidth=0.04cm](0.09734377,0.18796967)(4.4973435,0.18796967)
\psline[linewidth=0.04cm](2.2973437,2.2879696)(4.4973435,0.18796967)
\psline[linewidth=0.04cm](2.0973437,2.0879698)(2.2973437,1.8879696)
\psline[linewidth=0.04cm](2.2973437,1.8879696)(2.4973438,2.0879698)
\usefont{T1}{ptm}{m}{n}
\rput(0.54234374,0.37796965){$\theta$}
\usefont{T1}{ptm}{m}{n}
\rput(1.1764063,1.5979697){$a$}
% \rput (1, 2){1.}
\usefont{T1}{ptm}{m}{n}
\rput(2.1778126,-0.10203034){$b$}
\usefont{T1}{ptm}{m}{n}
\rput(3.759219,1.2979697){$c$}
\usefont{T1}{ptm}{m}{n}
\rput(6.596406,1.0179696){$a$}
\usefont{T1}{ptm}{m}{n}
\rput(7.479219,2.1179698){$c$}
\usefont{T1}{ptm}{m}{n}
\rput(8.797812,0.41796967){$b$}
\usefont{T1}{ptm}{m}{n}
\rput(13.076406,1.4979696){$a$}
\usefont{T1}{ptm}{m}{n}
\rput(10.777813,1.1979697){$b$}
\usefont{T1}{ptm}{m}{n}
\rput(12.659219,-0.20203033){$c$}
\usefont{T1}{ptm}{m}{n}
\rput(2.896407,-3.6220303){$a$}
\usefont{T1}{ptm}{m}{n}
\rput(2.2978134,-1.8220303){$b$}
\usefont{T1}{ptm}{m}{n}
\rput(4.2792187,-2.3220303){$c$}
\usefont{T1}{ptm}{m}{n}
\rput(6.596406,-1.8220303){$a$}
\usefont{T1}{ptm}{m}{n}
\rput(7.7978134,-3.1220303){$b$}
\usefont{T1}{ptm}{m}{n}
\rput(8.579218,-1.7220303){$c$}
\usefont{T1}{ptm}{m}{n}
\rput(10.496406,-2.6220303){$a$}
\usefont{T1}{ptm}{m}{n}
\rput(12.397812,-2.8220303){$b$}
\usefont{T1}{ptm}{m}{n}
\rput(11.579218,-1.4220303){$c$}
\rput{-66.69996}(-0.02258055,0.761629){\psarc[linewidth=0.04](0.5673438,0.39796966){0.35}{25.785578}{138.26501}}
\rput{106.83747}(12.386083,-6.856496){\psarc[linewidth=0.04](8.737344,1.1679693){0.34}{32.561234}{147.21713}}
\rput{93.56388}(15.110441,-13.562675){\psarc[linewidth=0.04](13.927344,0.31796986){0.39}{24.333387}{127.3505}}
\psarc[linewidth=0.04](14.697344,3.1279697){0.0}{0.0}{180.0}
\psarc[linewidth=0.04](14.837344,3.2279696){0.0}{0.0}{180.0}
\rput{153.95303}(6.6359906,-3.7989676){\psarc[linewidth=0.04](3.757346,-1.1320316){0.5}{50.00946}{137.84839}}
\psarc[linewidth=0.04](20.097343,3.9079697){0.0}{0.0}{180.0}
\psarc[linewidth=0.04](20.117344,3.8879697){0.0}{0.0}{180.0}
\psarc[linewidth=0.04](20.157345,3.8279696){0.0}{0.0}{180.0}
\rput{-46.143967}(4.180272,3.6097476){\psarc[linewidth=0.04](6.3273435,-3.1020286){0.57}{58.40874}{128.71255}}
\rput{-116.039185}(16.474606,6.6625834){\psarc[linewidth=0.04](10.317342,-1.8120308){0.52}{49.23777}{123.63383}}
\end{pspicture} 
}
\end{center}


\item Gebruik jou sakrekenaar om die waardes van die volgende te bepaal (korrek tot $2$ desimale plekke):
\begin{enumerate}[noitemsep, label=\textbf{(\alph*)} ]
\begin{multicols}{2} 
% \setcounter{enumi}{6} 
\item $tan~65^{\circ}$
\item $sin~38^{\circ}$
\item $cos~74^{\circ}$
\item $sin~12^{\circ}$
\item $cos~26^{\circ}$
\item $tan~49^{\circ}$
\item $\frac{1}{4}~cos~20^{\circ}$
\item $3~tan~40^{\circ}$
\item $\frac{2}{3}~sin~90^{\circ}$
\end{multicols}
\end{enumerate}

\item As $x=39^{\circ}$ en $y=21^{\circ}$ is, gebruik 'n sakrekenaar om te bepaal of die volgende waar of vals is:
\begin{enumerate}[noitemsep, label=\textbf{(\alph*)} ]
\begin{multicols}{2} 
% \setcounter{enumi}{15} 
\item $cos~x + 2~cos~x=3~cos~x$
\item $cos~2y = cos~y+cos~y$
\item $tan~x=\dfrac{sin~x}{cos~x}$
\item $cos~(x+y) = cos~x+cos~y$
\end{multicols}
\end{enumerate}


\item Voltooi elk van die volgende (die eerste voorbeeld is klaar gedoen):
\begin{center}
% \begin{minipage}{0.25 \textwidth}
\setcounter{subfigure}{0}
\scalebox{1}{
\begin{pspicture}(0,-2.0990624)(4.21625,2.0990624)
\psline[linewidth=0.04cm](0.4775,-1.5590625)(3.7775,1.5409375)
\psline[linewidth=0.04cm](3.7775,1.5409375)(3.7775,-1.5590625)
\psline[linewidth=0.04cm](0.4775,-1.5590625)(3.7775,-1.5590625)
\psline[linewidth=0.04cm](3.4775,-1.5590625)(3.4775,-1.2590625)
\psline[linewidth=0.04cm](3.4775,-1.2590625)(3.7775,-1.2590625)
\rput(4.025,1.8759375){$A$}
\rput(3.874375,-1.9240625){ $B$}
\rput(0.10625,-1.8240625){$C$}
\end{pspicture} 
}
% \end{minipage}
\end{center}

\begin{enumerate}[noitemsep, label=\textbf{(\alph*)} ]
\begin{multicols}{2}
% \setcounter{enumi}{19}
\item $sin~\hat{A} = \frac{\mbox{Teenoorstaande sy}}{\mbox{skuinssy}}=\dfrac{CB}{AC}$
\item $cos~\hat{A} = $
\item $tan~\hat{A}= $
\item $sin~\hat{C}= $
\item $cos~\hat{C}= $
\item $tan~\hat{C}= $
\end{multicols}
\end{enumerate}

\item Gebruik die driehoek hieronder om die volgende te voltooi:
\begin{center}
% \begin{minipage}{0.25\textwidth}
\scalebox{1} % Change this value to rescale die drawing.
{
\begin{pspicture}(0,-1.97)(5.3271875,1.97)
\psline[linewidth=0.04cm](0.361875,-1.95)(4.061875,-1.95)
\psline[linewidth=0.04cm](4.061875,-1.95)(4.061875,1.95)
\psline[linewidth=0.04cm](4.061875,1.95)(0.361875,-1.95)
\psline[linewidth=0.04cm](3.661875,-1.95)(3.661875,-1.55)
\psline[linewidth=0.04cm](3.661875,-1.55)(4.061875,-1.55)
\rput(1.7398437,0.26){$2$}
\rput(1.8,-2.2){$1$}
\rput(1.1554687,-1.64){$60^{\circ}$}
\rput(3.6554687,0.96){$30^{\circ}$}
\rput(4.405469,-0.44){$\sqrt{3}$}
\end{pspicture} 
}

% \end{minipage}
\end{center}


\begin{enumerate}[noitemsep, label=\textbf{(\alph*)} ]
\begin{multicols}{2}
% \setcounter{enumi}{25}
\item $sin~60^{\circ} = $
\item $cos~60^{\circ} = $
\item $tan~60^{\circ}= $
\item $sin~30^{\circ}= $
\item $cos~30^{\circ}= $
\item $tan~30^{\circ}= $
\end{multicols}
\end{enumerate}


\item Gebruik die driehoek hieronder om die volgende te voltooi:
\begin{center}
% \begin{minipage}{0.25\textwidth}
\scalebox{1} % Change this value to rescale die drawing.
{
\begin{pspicture}(0,-2.24)(5.0271873,2.22)
\psline[linewidth=0.04cm](0.361875,-1.7)(4.061875,-1.7)
\psline[linewidth=0.04cm](4.061875,-1.7)(4.061875,2.2)
\psline[linewidth=0.04cm](4.061875,2.2)(0.361875,-1.7)
\psline[linewidth=0.04cm](3.661875,-1.7)(3.661875,-1.3)
\psline[linewidth=0.04cm](3.661875,-1.3)(4.061875,-1.3)
\rput(2.2314062,-2.09){$1$}
\rput(1.1554687,-1.39){$45^{\circ}$}
\rput(3.7554688,1.41){$45^{\circ}$}
\rput(1.6054688,0.51){$\sqrt{2}$}
\rput(4.331406,0.21){$1$}
\end{pspicture} 
}
% \end{minipage}
\end{center}

\begin{enumerate}[noitemsep, label=\textbf{(\alph*)} ]

% \setcounter{enumi}{31}
\item $sin~45^{\circ} = $
\item $cos~45^{\circ} = $
\item $tan~45^{\circ}= $

\end{enumerate}
\end{enumerate}
}
\end{exercises}

\section{Spesiale hoeke}
Vir meeste hoeke $\theta $ het ons 'n sakrekenaar nodig om die waardes van $sin~\theta $, $cos~\theta $ en $tan~\theta $ te bereken. Ons het egter in die vorige oefening gesien dat ons hierdie waardes per hand kan bereken vir sekere spesiale hoeke. Die waardes van die trigonometriese funksies van hierdie hoeke word gelys in die tabel hieronder.\par
Onthou dat die lengtes van 'n reghoekige driehoek die Teorema van Pythagoras moet bevredig: die vierkant van die skuinssy is gelyk aan die som van die vierkante van die ander twee sye.\par 
% \textbf{m39408*id80733}\par
\begin{table}[H]
% \begin{table}[H]
% \\ '' '0'
\begin{center}
\setlength{\extrarowheight}{2.5pt}

\begin{tabular}{|m{1.2cm}|l|l|l|}\hline
&

${30}^{\circ }$
&
${45}^{\circ }$
&
${60}^{\circ }$


% make-rowspan-placeholders\tabularnewline\cline{1-1}\cline{2-2}\cline{3-3}\cline{4-4}\cline{5-5}\cline{6-6}\cline{7-7}

\\ \hline
%--------------------------------------------------------------------
$cos~\theta $
 &
$\frac{\sqrt{3}}{2}$
&
$\frac{1}{\sqrt{2}}$
&
$\frac{1}{2}$

\\ \hline
%--------------------------------------------------------------------
$sin~\theta $
&

$\frac{1}{2} $
&
$\frac{1}{\sqrt{2}}$
&
$\frac{\sqrt{3}}{2}$

%  make-rowspan-placeholders
\\ \hline
%--------------------------------------------------------------------
$tan~\theta $
&

$\frac{1}{\sqrt{3}}$
&
$1$ &
$\sqrt{3}$



% make-rowspan-placeholders
\\ \hline
%--------------------------------------------------------------------
\end{tabular}
\end{center}
% \begin{center}{\small\bfseries Table 14.3}\end{center}
% \begin{caption}{\small\bfseries Table 14.3}\end{caption}
\end{table}


\begin{minipage}{0.5\textwidth}
\begin{center}
\scalebox{0.7} % Change this value to rescale die drawing.
{
\begin{pspicture}(0,-1.97)(5.3271875,1.97)
\psline[linewidth=0.04cm](0.361875,-1.95)(4.061875,-1.95)
\psline[linewidth=0.04cm](4.061875,-1.95)(4.061875,1.95)
\psline[linewidth=0.04cm](4.061875,1.95)(0.361875,-1.95)
\psline[linewidth=0.04cm](3.661875,-1.95)(3.661875,-1.55)
\psline[linewidth=0.04cm](3.661875,-1.55)(4.061875,-1.55)
\rput(1.7398437,0.26){\LARGE$2$}
\rput(1.8,-2.2){\LARGE$1$}
\rput(1.1554687,-1.64){\LARGE$60^{\circ}$}
\rput(3.6554687,0.96){\LARGE$30^{\circ}$}
\rput(4.405469,-0.44){\LARGE$\sqrt{3}$}
\end{pspicture} 
}
\end{center}
\end{minipage}


\begin{minipage}{0.5\textwidth}
\begin{center}
\scalebox{0.7} % Change this value to rescale die drawing.
{
\begin{pspicture}(0,-2.24)(5.0271873,2.22)
\psline[linewidth=0.04cm](0.361875,-1.7)(4.061875,-1.7)
\psline[linewidth=0.04cm](4.061875,-1.7)(4.061875,2.2)
\psline[linewidth=0.04cm](4.061875,2.2)(0.361875,-1.7)
\psline[linewidth=0.04cm](3.661875,-1.7)(3.661875,-1.3)
\psline[linewidth=0.04cm](3.661875,-1.3)(4.061875,-1.3)
\rput(2.2314062,-2.09){\LARGE$1$}
\rput(1.1554687,-1.39){\LARGE$45^{\circ}$}
\rput(3.7554688,1.41){\LARGE$45^{\circ}$}
\rput(1.6054688,0.51){\LARGE$\sqrt{2}$}
\rput(4.331406,0.21){\LARGE$1$}
\end{pspicture} 
}
\end{center}
\end{minipage}

\par
Hierdie waardes is handig wanneer ons 'n probleem in trigonometrie moet oplos sonder die gebruik van 'n sakrekenaar.\par 
\begin{exercises}{}
{
\begin{enumerate}[itemsep=6pt, label=\textbf{\arabic*}. ] 
\item Bereken die volgende sonder 'n sakrekenaar:
\begin{enumerate}[noitemsep, label=\textbf{(\alph*)} ]
\item $sin~45^{\circ} \times cos~45^{\circ}$
\item $cos~60^{\circ} + tan~45^{\circ}$
\item $sin~60^{\circ} - cos~60^{\circ}$
\end{enumerate}

\item gebruik die volgende tabel om aan te dui dat:
\begin{enumerate}[itemsep=5pt, label=\textbf{(\alph*)} ]
% \setcounter{enumi}{3}
\item $\dfrac{sin~60^{\circ}}{cos~60^{\circ}} = tan~60^{\circ} $
\item $sin^{2}~45^{\circ}+ cos^{2}~45^{\circ} =1 $
\item $cos~30^{\circ} =\sqrt{1- sin^{2}}~30^{\circ}$
\end{enumerate}

\item Gebruik die definisies van die trigonometriese verhoudings om die volgende vrae te antwoord:
\begin{enumerate}[noitemsep, label=\textbf{(\alph*)} ]
% \setcounter{enumi}{6}
\item Verduidelik hoekom $sin~\alpha \leq 1$ is vir alle $\alpha$.
\item Verduidelik hoekom $cos~\beta$ 'n maksimum waarde van $1$ het.
\item Is daar 'n maksimum vir $tan~\gamma$ ?
\end{enumerate}
\end{enumerate}
}
\end{exercises}


\section{Omgekeerde funksies}
Elk van die drie trigonometriese funksies het 'n omgekeerde. Die omgekeerdes cosecant, secant and cotangens word soos volg gedefinieer:

\begin{equation*}
\begin{array}{ccc}cosec~\theta & =& \dfrac{1}{sin~\theta } \vspace{3pt}\\
 sec~\theta & =& \dfrac{1}{cos~\theta } \vspace{3pt}\\
 cot~\theta & =& \dfrac{1}{tan~\theta }
\end{array}
\end{equation*}
Ons kan ook hierdie omgekeerdes definieer vir enige reghoekige driehoek:

\begin{equation*}
\begin{array}{ccc}\hfill cosec~\theta & =& \frac{\mbox{skuinssy}}{\mbox{teenoorstaande sy}}\hfill \vspace{3pt}\\
 \hfill sec~\theta & =& \frac{\mbox{skuinssy}}{\mbox{aanliggende sy}}\hfill \vspace{3pt}\\
 \hfill cot~\theta & =& \frac{\mbox{aanliggende sy}}{\mbox{teenoorstaande sy}}\hfill 
\end{array}
\end{equation*}
Note that:
\begin{equation*}
\begin{array}{ccc}sin~\theta \times cosec~\theta & =& 1 \vspace{3pt}\\
 cos~\theta \times sec~\theta & =& 1 \vspace{3pt}\\
tan~\theta \times cot~\theta & =& 1
\end{array}
\end{equation*}
\section{Oplossing van trigonometriese vergelykings}

\begin{wex}{Berekening van lengtes}{Bereken die lengte van $x$ in die volgende reghoekige driehoek: \\
\begin{center}
\scalebox{1} 
{
\begin{pspicture}(0,-1.02)(4.655469,2.02)
\psline[linewidth=0.04](4.0,-1.0)(1.0,-1.0)(4.0,2.0)(4.0,-1.0)
\rput(2.0275,0.77){$100$}
\psarc[linewidth=0.04](1.0,-1.0){1.0}{0.0}{45.0}
\rput(1.6450001,-0.75){$50^{\circ}$}
\rput(4.355,0.41){$x$}
\psline[linewidth=0.04cm](3.6,-0.98)(3.6,-0.58)
\psline[linewidth=0.04cm](3.6,-0.58)(4.0,-0.58)
\end{pspicture} 
}
\end{center}
}
{
\westep{Identifiseer die teenoorstaande sy en die aanliggende sy}
\begin{equation*}
\begin{array}{ccl}
 
\hfill sin~\theta &=& \dfrac{\mbox{teenoorstaande sy}}{\mbox{skuinssy}}  \hfill \vspace{5pt}\\
\hfill sin~ 50^\circ &=& \dfrac{x}{100}  \hfill \\
\end{array}
\end{equation*}



\westep{Herrangskik die vergelyking om vir $x$ op te los}
\begin{equation*}
 x=100 \times sin~50^{\circ}
\end{equation*}

\westep{Gebruik jou sakrekenaar om die antwoord te bereken}
\begin{equation*}
x = 76.6
\end{equation*}
}
\end{wex}




Die volgende videos gee 'n opsomming van wat jy tot nou toe geleer het.
\setcounter{subfigure}{0}
\begin{figure}[H] % horizontal\label{m39408*trigonometry-1}
\textnormal{Khan Academy video on trigonometry - 1}\vspace{.1in} \nopagebreak
\label{m39408*yt-media}\label{m39408*yt-video}
\raisebox{-5 pt}{ \includegraphics[width=0.5cm]{col11306.imgs/summary_www.png}} { (Video:  MG10102 )}
\vspace{2pt}
\vspace{.1in}
\end{figure}       \par \label{m39408*eip-33}
\setcounter{subfigure}{0}
\begin{figure}[H] % horizontal\label{m39408*trigonometry-2}
\textnormal{Khan Academy video on trigonometry - 2}\vspace{.1in} \nopagebreak
\label{m39408*yt-media2}\label{m39408*yt-video2}
\raisebox{-5 pt}{ \includegraphics[width=0.5cm]{col11306.imgs/summary_www.png}} { (Video:  MG10103 )}
\vspace{2pt}
\vspace{.1in}
\end{figure}       


\begin{exercises}{}
{
\begin{enumerate}[itemsep=5pt, label=\textbf{\arabic*}. ]
\item Vir elke driehoek, bereken die lengte van die sy met 'n letter gemerk. Gee die antwoorde korrek tot $2$ desimale plekke.
\begin{center}
\scalebox{0.85} % Change this value to rescale die drawing.
{
\begin{pspicture}(0,-4.11)(11.1198435,5)
\psline[linewidth=0.04](4.974369,2.612122)(2.6350498,3.9603012)(0.6377475,0.49464345)(4.974369,2.612122)
\psline[linewidth=0.04cm](2.8949742,3.8105037)(2.7451768,3.550579)
\psline[linewidth=0.04cm](2.7451768,3.550579)(2.4852521,3.7003772)
\psline[linewidth=0.04](4.557088,-3.2752385)(4.589271,-0.5754303)(0.5895552,-0.52775127)(4.557088,-3.2752385)
\psline[linewidth=0.04cm](4.5856953,-0.8754088)(4.2857165,-0.87183315)
\psline[linewidth=0.04cm](4.2857165,-0.87183315)(4.2892923,-0.5718545)
\psline[linewidth=0.04](7.0587263,3.3569472)(6.922255,0.6603981)(10.917143,0.45821914)(7.0587263,3.3569472)
\psline[linewidth=0.04cm](6.9374194,0.9600146)(7.2370353,0.94485116)
\psline[linewidth=0.04cm](7.2370353,0.94485116)(7.2218714,0.64523476)
\psline[linewidth=0.04](7.9354,-3.9493167)(10.598035,-3.9603012)(10.613694,-0.16448934)(7.9354,-3.9493167)
\psline[linewidth=0.04cm](10.340361,-3.959238)(10.341382,-3.7116852)
\psline[linewidth=0.04cm](10.341382,-3.7116852)(10.599056,-3.7127483)
% \usefont{T1}{ptm}{m}{n}
\rput(0.12453125,3.7196066){\textbf{(a)}}
% \usefont{T1}{ptm}{m}{n}
\rput(6.133125,3.7396066){\textbf{(b)}}
% \usefont{T1}{ptm}{m}{n}
\rput(0.11421875,-0.26039344){\textbf{(c)}}
% \usefont{T1}{ptm}{m}{n}
\rput(6.1343746,-0.28039345){\textbf{(d)}}
\usefont{T1}{ptm}{m}{n}
\rput(4.197969,2.6346066){$37^\circ$}
\usefont{T1}{ptm}{m}{n}
\rput(3.0853126,1.2){$62$}
\usefont{T1}{ptm}{m}{n}
\rput(1.3824998,2.5846066){$a$}
\usefont{T1}{ptm}{m}{n}
\rput(9.804531,0.8146064){$23^\circ$}
\usefont{T1}{ptm}{m}{n}
\rput(8.595937,0.21460645){$21$}
\usefont{T1}{ptm}{m}{n}
\rput(6.663907,2.0646067){$b$}
\usefont{T1}{ptm}{m}{n}
\rput(4.2134376,-2.6453934){$ 55^\circ$}

\usefont{T1}{ptm}{m}{n}
\rput(2.166719,-2.1653934){$19$}
\usefont{T1}{ptm}{m}{n}
\rput(4.8759375,-1.8353934){$c$}
\usefont{T1}{ptm}{m}{n}
\rput(8.953907,-1.9653934){$33$}
\usefont{T1}{ptm}{m}{n}
\rput(10.22,-1.1453935){ $49^\circ$}
\usefont{T1}{ptm}{m}{n}
\rput(10.90422,-2.0153935){$d$}
\end{pspicture}    
}
\end{center}




\begin{center}
 \scalebox{0.85} % Change this value to rescale die drawing.
{
\begin{pspicture}(0,-3.8416457)(11.207031,3.8216705)
\psline[linewidth=0.04](2.3197074,3.8016455)(1.0997375,1.3929795)(4.668132,-0.4143837)(2.3197074,3.8016455)
\psline[linewidth=0.04cm](1.2352896,1.6606089)(1.5029193,1.5250567)
\psline[linewidth=0.04cm](1.5029193,1.5250567)(1.3673671,1.2574271)
% \usefont{T1}{ptm}{m}{n}
\rput(0.11234375,3.5211668){\textbf{(e)}}
% \usefont{T1}{ptm}{m}{n}
\rput(6.084844,3.5211668){\textbf{(f)}}
% \usefont{T1}{ptm}{m}{n}
\rput(0.14375,-0.39883313){\textbf{(g)}}
% \usefont{T1}{ptm}{m}{n}
\rput(6.1403127,-0.498833){\textbf{(h)}}
\psline[linewidth=0.04](1.0870312,-0.6888331)(1.0870312,-3.3888335)(5.0870314,-3.3888335)(1.0870312,-0.6888331)
\psline[linewidth=0.04cm](1.0870312,-3.0888333)(1.3870313,-3.0888333)
\psline[linewidth=0.04cm](1.3870313,-3.0888333)(1.3870313,-3.3888335)
\psline[linewidth=0.04](10.705838,-3.6336784)(10.738021,-0.93386954)(6.738305,-0.88619083)(10.705838,-3.6336784)
\psline[linewidth=0.04cm](10.734446,-1.2338486)(10.434467,-1.2302725)
\psline[linewidth=0.04cm](10.434467,-1.2302725)(10.438043,-0.9302939)
\psline[linewidth=0.04](11.200404,0.87534267)(9.325028,2.8177538)(6.4473815,0.039416883)(11.200404,0.87534267)
\psline[linewidth=0.04cm](9.533403,2.6019304)(9.31758,2.393555)
\psline[linewidth=0.04cm](9.31758,2.393555)(9.109204,2.6093786)
\usefont{T1}{ptm}{m}{n}
\rput(3.5887504,1.9261669){$e$}
\usefont{T1}{ptm}{m}{n}
\rput(3.8853126,0.3161668){$17^\circ$}
\usefont{T1}{ptm}{m}{n}
\rput(1.3671876,2.6761668){ $12$}
\usefont{T1}{ptm}{m}{n}
\rput(7.533282,0.51616687){$22^\circ$}
\usefont{T1}{ptm}{m}{n}
\rput(9.139219,0.2861668){$f$}
\usefont{T1}{ptm}{m}{n}
\rput(7.758125,1.7361668){ $31$}
\usefont{T1}{ptm}{m}{n}
\rput(3.028594,-1.683833){$32$}
\usefont{T1}{ptm}{m}{n}
\rput(4.1132817,-3.123833){$23^\circ$}
\usefont{T1}{ptm}{m}{n}
\rput(2.949844,-3.6138332){$g$}
\usefont{T1}{ptm}{m}{n}
\rput(10.928438,-2.2538335){$h$}
\usefont{T1}{ptm}{m}{n}
\rput(7.7667193,-1.2438331){ $30^\circ$}
\usefont{T1}{ptm}{m}{n}
\rput(8.245001,-2.3038335){$20$}
\end{pspicture} 
}
   
\end{center}  

\item Skryf twee verhoudings neer in terme van die sye vir elk van die volgende: $AB$; $BC$; $BD$; $AD$; $DC$ and $AC$\\
\begin{center}
\scalebox{1} % Change this value to rescale die drawing.
{
\begin{pspicture}(0,-1.2515883)(4.02125,1.2515885)
\psline[linewidth=0.04,fillstyle=solid](0.019002499,-1.2315885)(0.009700612,1.2315885)(4.00125,-1.1902716)(0.039492033,-1.2120903)(0.039492033,-1.2120903)(0.0,-1.1911051)
\psline[linewidth=0.04](0.009700612,-0.9884115)(0.20970061,-0.9884115)(0.20970061,-1.2284116)
\psline[linewidth=0.04,fillstyle=solid](1.3697007,0.4315885)(0.049700614,-1.1684115)(0.049700614,-1.1484115)
\psline[linewidth=0.04](1.2297006,0.2715885)(1.0697006,0.3515885)(1.1897006,0.5115885)
\usefont{T1}{ptm}{m}{n}
\rput(0,-1.5){$A$}
\rput(0, 1.5){$B$}
\rput(1.4, 0.7){$C$}
\rput(4, -1.5){$D$}
\end{pspicture} 
}
\end{center}
     \begin{enumerate}[noitemsep, label=\textbf{(\alph*)} ]
    \item $sin~\hat{B}$
    \item $cos~\hat{D}$
    \item $tan~\hat{B}$
    \end{enumerate}
\vspace{10pt}
\item In $\triangle MNP$, $\hat{N}=90^{\circ}$, $MP=20$ en $\hat{P}=40^{\circ}$. Bereken $NP$ en $MN$ (korrek tot $2$ desimale plekke).
\end{enumerate}

\par 

\par \raisebox{-0.2em}{\includegraphics[height=1em]{../icons/www.pdf}} Vind die antwoorde met die kort kodes:
\par \begin{tabular}[h]{cccccc}
(1.) lc1  & \end{tabular}
}
\end{exercises}

\section{Berekening van 'n hoek}
 
\begin{wex}{Hoekberekeninge}
{
Vind die waarde van $\theta$ in die volgende reghoekige driehoek: \\
\begin{center}
\scalebox{1}  
{ 
\begin{pspicture}(0,-1.3928125)(5.0089064,2.02)
\psline[linewidth=0.04](4.0,-1.0)(1.0,-1.0)(4.0,2.0)(4.0,-1.0)
\usefont{T1}{ptm}{m}{n}
\rput(2.7292187,-1.19){$100$}
\usefont{T1}{ptm}{m}{n}
\rput(4.5243754,0.61){$50$}
\usefont{T1}{ptm}{m}{n}
\rput(1.5981251,-0.69){$\theta$}
\psarc[linewidth=0.04](1.0,-1.0){1.0}{0.0}{44.06081}
\psline[linewidth=0.04](4.02,-0.7)(3.68,-0.7)(3.68,-0.98)(3.68,-0.98)
\end{pspicture} 
 
}
\end{center}
}
{

\westep{Identifiseer die teenoorstaande- en aanliggende sye en die skuinssy}
In hierdie geval het jy die teenoorstaande sy en die aanliggende vir hoek $\theta$. \\

\begin{equation*}
\begin{array}{ccl}
 
\hfill tan~\theta &=& \dfrac{\mbox{teenoorstaande}}{\mbox{aanliggend}}  \hfill \vspace{5pt}\\
\hfill tan~ \theta &=& \dfrac{50}{100} \hfill \\
\end{array}
\end{equation*}
% There is a horrible error here: The example show opp and hypotenuse is given, but die calc is for opp and adjacent. In die proper
% English version as well. I am correcting this to read opp and adjacent.
\westep{Gebruik jou sakrekenaar om op te los vir $\theta$ }
Om op te los vir $\theta$ sal jy die inverse funksie van jou sakrekenaar benodig: \vspace{10pt}
\\
Druk \fbox{2ndF} \fbox{TAN} \fbox{(} \fbox{50} \fbox{\div} \fbox{100} \fbox{)} \fbox{\LARGE =} $26,6$
\westep{Skryf die finale antwoord neer}
\begin{equation*}
\theta = 26,6^{\circ}
\end{equation*}

}
\end{wex}

\begin{exercises}{}
{
   \begin{enumerate}[itemsep=5pt, label=\textbf{\arabic*}. ] 
\item Bepaal die hoek (korrek tot $1$ desimale plek):
    \begin{enumerate}[itemsep=3pt, label=\textbf{(\alph*)} ]
\begin{multicols}{2}
 \item $tan~\theta = 1,7$
\item $sin~\theta = 0,8$
\item $cos~\alpha = 0,32$
\item $tan~\theta = 5\frac{3}{4}$
\item $sin~\theta = \frac{2}{3}$
\item $cos~\gamma = 1,2$
\item $4~cos~\theta = 3$
\item $cos~4\theta = 0,3$
\item $sin~\beta + 2= 2,65$
\item $sin~\theta = 0,8$
\item $3 ~tan~\beta = 1$
\item $sin~3\alpha = 1,2$
\item $tan~(\frac{\theta}{3}) = sin~48^{\circ}$
\item $\frac{1}{2}~cos~2\beta = 0,3$
\item $2~sin~3\theta +1= 2,6$
\end{multicols}
\end{enumerate}

\item Bepaal $\alpha$ in die volgende reghoekige driehoeke:
\begin{center}
\scalebox{1} % Change this value to rescale die drawing.
{
\begin{pspicture}(0,-4.032389)(7.3778124,4.17849)

\rput(0, 3.5){\textbf{(a)}}
\rput(3.5, 3.5){\textbf{(b)}}
\rput(0, 0.5){\textbf{(c)}}
\rput(3.5, 0.5){\textbf{(d)}}
\rput(0, -2.0){\textbf{(e)}}
\rput(3.5, -2.0){\textbf{(f)}}
\psline[linewidth=0.04,fillstyle=solid](4.128174,1.9359901)(4.129576,3.691884)(6.1603127,1.9559901)(4.142015,1.9498184)(4.142015,1.9498184)(4.115519,1.9649143)
\psline[linewidth=0.04,fillstyle=solid](0.5403125,3.7206802)(2.3103564,3.7223134)(0.5605087,1.1162983)(0.5542749,3.7072453)(0.5542749,3.7072453)(0.5694489,3.7330344)
\psline[linewidth=0.04](0.8203125,3.71599)(0.8203125,3.41599)(0.5403125,3.41599)
\psline[linewidth=0.04](4.132775,2.2306225)(4.429114,2.2301245)(4.4286456,1.9555349)
\psline[linewidth=0.04,fillstyle=solid](0.58817375,-1.2040099)(0.5895763,0.5518841)(2.6203125,-1.1840099)(0.60201496,-1.1901817)(0.60201496,-1.1901817)(0.57551914,-1.1750857)
\psline[linewidth=0.04](0.59277505,-0.90937746)(0.889114,-0.9098756)(0.88864577,-1.184465)
\psline[linewidth=0.04,fillstyle=solid](0.57534015,-1.9803641)(2.331222,-1.9737015)(0.6046739,-4.0123897)(0.5892318,-1.9941416)(0.5892318,-1.9941416)(0.6042059,-1.9675767)
\psline[linewidth=0.04](0.8699906,-1.983612)(0.8708536,-2.2799501)(0.5962649,-2.2807431)
\psline[linewidth=0.04,fillstyle=solid](6.6596694,0.2178756)(6.6803126,-0.6840099)(3.9803126,0.21599011)(6.645639,0.20423959)(6.645639,0.20423959)(6.6719236,0.18877967)
\psline[linewidth=0.04](6.6710052,-0.076665334)(6.3803124,-0.06400988)(6.3789563,0.20247632)
\psline[linewidth=0.04,fillstyle=solid](5.604427,-1.7804683)(6.7567835,-3.1053228)(4.0853443,-3.1304228)(5.603074,-1.7999867)(5.603074,-1.7999867)(5.632968,-1.7939644)
\psline[linewidth=0.04](5.794496,-2.005642)(5.570731,-2.199926)(5.3907113,-1.9925796)
\rput{-34.695152}(-0.9369147,0.7926677){\psarc[linewidth=0.024](0.8003125,1.8959901){0.3}{39.289406}{180.0}}
\usefont{T1}{ptm}{m}{n}
\rput(1.4359375,3.9229398){$4$}
\usefont{T1}{ptm}{m}{n}
\rput(0.29078126,2.4629397){$9$}
\usefont{T1}{ptm}{m}{n}
\rput(5.470625,3.02294){$13$}
\usefont{T1}{ptm}{m}{n}
\rput(5.034219,1.5829399){$7,5$}
\usefont{T1}{ptm}{m}{n}
\rput(1.75375,-0.057060193){$2,2$}
\usefont{T1}{ptm}{m}{n}
\rput(0.2096875,-0.4770602){$1,7$}
\usefont{T1}{ptm}{m}{n}
\rput(5.716875,0.4829398){$9,1$}
\usefont{T1}{ptm}{m}{n}
\rput(7.0603123,-0.3170602){$4,5$}
\usefont{T1}{ptm}{m}{n}
\rput(0.26703125,-2.8770602){$12$}
\usefont{T1}{ptm}{m}{n}
\rput(1.9092188,-3.1370602){$15$}
\usefont{T1}{ptm}{m}{n}
\rput(4.5784373,-2.3170602){$1$}
\usefont{T1}{ptm}{m}{n}
\rput(5.54375,-3.47706){$\sqrt{2}$}
\usefont{T1}{ptm}{m}{n}
\rput(6.2634373,-2.93706){$\alpha$}
\usefont{T1}{ptm}{m}{n}
\rput(0.7834375,-3.4970603){$\alpha$}
\usefont{T1}{ptm}{m}{n}
\rput(1.9234375,-0.9970602){$\alpha$}
\usefont{T1}{ptm}{m}{n}
\rput(5.1634374,-0.017060194){$\alpha$}
\usefont{T1}{ptm}{m}{n}
\rput(4.3034377,3.1429398){$\alpha$}
\usefont{T1}{ptm}{m}{n}
\rput(0.7634375,1.8029398){$\alpha$}
\rput{-162.53017}(7.4652467,7.66313){\psarc[linewidth=0.024](4.3213253,3.258065){0.34310362}{39.289406}{156.30495}}
\rput{-89.91282}(5.095528,5.0594177){\psarc[linewidth=0.024](5.081325,-0.02193504){0.34310362}{53.61859}{134.10115}}
\rput{-319.68652}(-0.17349519,-1.5915118){\psarc[linewidth=0.024](2.0811038,-1.0320795){0.42839816}{67.16096}{156.30495}}
\rput{-30.447884}(1.9160283,-0.06363264){\psarc[linewidth=0.024](0.8411037,-3.5520794){0.42839816}{67.16096}{156.30495}}
\rput{-319.68652}(-0.41087502,-4.8648257){\psarc[linewidth=0.024](6.4211035,-2.9920795){0.42839816}{67.16096}{156.30495}}
\end{pspicture} 

}
\end{center}
\end{enumerate}
}
\end{exercises}



\section{Twee-dimensionele probleme}
Trigonometrie is ontwikkel deur oerbeskawings om praktiese probleme op te los in konstruksie en navigasie. Ons sal wys dat trigonometrie ook gebruik kan word om sekere ander praktiese probleme op te los. Ons gebruik trigonometriese funksies om probleme in twee dimensies op te los met behulp van reghoekige driehoeke.

\begin{wex}{Die vlieg van 'n vlieer}
{Mandla vlieg 'n vlieer aan 'n $17~$m lyn teen 'n skuinste van $63^{\circ}$.
\begin{enumerate}[noitemsep, label=\textbf{\arabic*}. ] 
 \item Wat is die hoogte ($h$) van die vlieer bo die grond?
\item As Mandla se vriend Sipho direk onder die vlieer staan, bereken die afstand ($d$) tussen die twee vriende. 
\end{enumerate}
}
{
\westep{Teken 'n skets en identifiseer die teenoorstaande sy, die aanliggende sy en die skuinssy}
\begin{center}
\scalebox{1} % Change this value to rescale die drawing.
{
\scalebox{1} % Change this value to rescale die drawing.
{
\begin{pspicture}(0,-1.9664062)(4.6114063,1.9364063)
\psline[linewidth=0.025999999](4.04,1.3864063)(4.04,1.4664062)(4.02,-1.4735937)(0.36,-1.4735937)(4.04,1.4464062)
\usefont{T1}{ppl}{m}{n}
\rput(4.236875,-0.12359375){$h$}
\usefont{T1}{ppl}{m}{n}
\rput(2.3425002,-1.7635938){$d$}
\usefont{T1}{ppl}{m}{n}
\rput(1.0521874,-1.2435937){$63$}
\pscircle[linewidth=0.02,dimen=outer](4.28,-0.80728817){0.1}
\psline[linewidth=0.02cm](4.28,-0.9072884)(4.28,-1.2935938)
\psline[linewidth=0.02cm](4.14,-1.4535937)(4.28,-1.2935938)
\psline[linewidth=0.02cm](4.28,-1.2881135)(4.4,-1.4535937)
\psline[linewidth=0.02cm](4.16,-1.0760568)(4.4,-1.0760568)
\usefont{T1}{ppl}{m}{n}
\rput(2.2839062,0.41640624){$17$}
\psline[linewidth=0.02](3.96,1.7864063)(4.12,1.9264063)(4.22,1.7597395)(4.0466666,1.4664062)(3.9773333,1.7760359)(4.22,1.7597395)
\psline[linewidth=0.02cm](4.04,1.4864062)(4.12,1.9064063)
\pscircle[linewidth=0.02,dimen=outer](0.14,-0.82728827){0.1}
\psline[linewidth=0.02cm](0.14,-0.92728823)(0.14,-1.3135937)
\psline[linewidth=0.02cm](0.0,-1.4735937)(0.14,-1.3135937)
\psline[linewidth=0.02cm](0.14,-1.3081135)(0.26,-1.4735937)
\psline[linewidth=0.02cm](0.02,-1.0960567)(0.3,-1.0935937)
\psline[linewidth=0.04](4.02,-1.2135937)(3.74,-1.2135937)(3.74,-1.4735937)
\end{pspicture} 
}
}
\end{center}

\westep{Gebruik die gegewens en toepaslike verhoudings om op te los vir $h$ and $d$}
\begin{enumerate}[noitemsep, label=\textbf{\arabic*}. ] 
\item
\begin{equation*}
 \begin{array}{ccl}
\hfill sin~63^{\circ} &= & \frac{\mbox{\footnotesize teenoorstaande}}{\mbox{\footnotesize skuinssy}} \hfill \vspace{5pt}\\
\hfill sin~63^{\circ} &=& \dfrac{h}{17} \hfill \\
\hfill \therefore h &= & 17~sin ~63^{\circ} \hfill \\
\hfill &=& 15,15~\mbox{m} \hfill \\
   \end{array}
\end{equation*}

\item
\begin{equation*}
 \begin{array}{ccl}
\hfill cos~63^{\circ} &= & \frac{\mbox{\footnotesize aanliggend}}{\mbox{\footnotesize skuinssy}} \hfill \vspace{5pt}\\
\hfill cos~63^{\circ} &=& \dfrac{d}{17} \hfill \\
\hfill \therefore d &= & 17~cos ~63^{\circ} \hfill \\
\hfill  &=& 7,72~\mbox{m} \hfill \\
\end{array}
\end{equation*}
\end{enumerate}
Let op dat die derde sy van die driehoek ook bereken kan word deur die Teorema van Pythagoras: $d^{2} = 17^{2} - h^{2}$.

\westep{Skryf die finale antwoord neer}
\begin{enumerate}[noitemsep, label=\textbf{\arabic*}. ] 
\item Die vlieer is $15,15$ m bo die grond.
\item Mandla and Sipho is $7,72$ m van mekaar.
\end{enumerate}
}
 
\end{wex}

\begin{wex}
{Berekening van hoeke}
{
ABCD is 'n trapesium met $AB=4\mbox{cm}$, $CD=6\mbox{cm}$, $BC=5\mbox{cm}$ and $AD=5\mbox{cm}$. Punt $E$ op diagonaal $AC$ verdeel die diagonaal so dat $AE=3\mbox{cm}$. $B\hat{E}C = 90^{\circ}$. Vind $A\hat{B}C$.
}
{
\westep{Teken die trapesium en merk al die gegewe lengtes op die diagram. Dui aan dat $B\hat{E}C$ 'n regte hoek is}
\begin{center}
\scalebox{1} % Change this value to rescale die drawing.
{
\begin{pspicture}(0,-2.8215597)(7.3365626,2.6459403)
\psline[linewidth=0.04](7.0446873,-2.3765597)(0.2846875,-2.3565598)(1.9846874,2.1834402)(5.5246873,2.1834402)(7.0246873,-2.3565598)
\psline[linewidth=0.04cm](2.0246875,2.1634402)(7.0246873,-2.3765597)
\psline[linewidth=0.04cm,linestyle=dashed,dash=0.16cm 0.16cm](5.5046873,2.1634402)(3.9646876,0.44344023)
\psline[linewidth=0.04](4.1846876,0.6834402)(4.4246874,0.48344022)(4.1646876,0.20344023)
\rput{124.910225}(10.204001,-1.1108128){\psarc[linewidth=0.04](5.3916802,2.1056097){0.48547706}{48.07423}{180.0}}
% \usefont{T1}{ptm}{m}{n}
\rput(5.7178125,2.4534402){$B$}
% \usefont{T1}{ptm}{m}{n}
\rput(7.23,-2.2465599){$C$}
% \usefont{T1}{ptm}{m}{n}
\rput(1.885,2.4534402){$A$}
% \usefont{T1}{ptm}{m}{n}
\rput(0.075,-2.2865598){$D$}
% \usefont{T1}{ptm}{m}{n}
\rput(3.7692187,0.15344024){$E$}
% \usefont{T1}{ptm}{m}{n}
\rput(3.7959375,2.4134402){$4$ cm}
% \usefont{T1}{ptm}{m}{n}
\rput(0.61078125,0.27344024){$5$ cm}
% \usefont{T1}{ptm}{m}{n}
\rput(6.790781,0.19344023){$5$ cm}
% \usefont{T1}{ptm}{m}{n}
\rput{-43.673504}(0.04067405,2.2905717){\rput(2.8509376,1.0734402){$3$ cm}}
% \usefont{T1}{ptm}{m}{n}
\rput(3.4921875,-2.6665597){$6$ cm}
\end{pspicture} 
}
\end{center}
     

\westep{Gebruik $\triangle ABE$ and $\triangle CBE$ om die twee hoeke by $\hat{B}$ te bepaal} 

\westep{Vind die eerst hoek $A\hat{B}E$}
Die Skuinssy en teenoorstaande sy word gegee vir beide driehoeke, dus gebruik ons die $sin$ funksie:
In $\triangle ABE$, \\

\begin{equation*}
\begin{array}{ccl}\hfill sin~(A\hat{B}E)& =& \frac{\mbox{\footnotesize teenoorstaande sy}}{\mbox{\footnotesize Skuinssy}}\hfill \vspace{5pt}\\
 \hfill & =& \dfrac{3}{4}\hfill \\
 \therefore A\hat{B}E& =& {48,59}^{\circ }\hfill 
\end{array}
\end{equation*}

\westep{Gebruik die Teorema van Pythagoras om $EC$ te bepaal}
In $\triangle ABE$, \\
\begin{equation*}
\begin{array}{ccl}
 \hfill BE^{2} &=& AB^{2} - AE^{2} \hfill \\
\hfill &=& 4^{2} - 3^{2} \hfill \\
\hfill &=&7\\
\therefore BE &=& \sqrt{7}
\end{array}
\end{equation*}

In $\triangle CBE$, \\

\begin{equation*}
\begin{array}{ccl}
 \hfill EC^{2} &=& BC^{2} - BE^{2} \hfill \\
\hfill &=& 5^{2} - (\sqrt{7})^{2} \hfill \\
\hfill &=& 18 \hfill \\
\therefore EC &=& 4,24 \hfill
\end{array}
\end{equation*}

\westep{Vind die tweede hoek $C\hat{B}E$}
In $\triangle CBE$ 
\begin{equation*}
\begin{array}{ccl}\hfill sin~(C\hat{B}E)& =& \frac{\mbox{\footnotesize teenoorstaande sy}}{\mbox{\footnotesize skuinssy}}\hfill \vspace{5pt}\\
 \hfill & =& \dfrac{4,24}{5}\hfill \\
 \therefore C\hat{B}E& =& {58}^{\circ }\hfill 
\end{array}
\end{equation*}

\westep{Bereken die som van die hoeke}
\begin{equation*}
A\hat{B}C = 48,59^{\circ} + 58^{\circ} = 106,59^{\circ}
\end{equation*}
}
\end{wex}



Nog 'n toepassing is om trigonometrie te gebruik om die hoogte van 'n gebou te bepaal. Ons kan 'n maatband van die dak laat hang, maar dit is onprakties (en gevaarlik) vir hoe geboue. Dis veel meer sinvol om trigonometrie te gebruik.\par 

\begin{wex}{Bepaling van 'n gebou se hoogte}
{
% \setcounter{subfigure}{0}
% \begin{figure}[htbp]
\begin{center}
\scalebox{1} % Change this value to rescale die drawing.
{
\begin{pspicture}(0,-2.5182812)(7.84,2.5717187)
\psframe[linewidth=0.028222222,dimen=outer,fillstyle=crosshatch,hatchwidth=0.028222222,hatchangle=180.0](7.84,2.0682812)(5.84,-1.9317187)
\psline[linewidth=0.028222222cm](0.84,-1.9317187)(5.84,-1.9317187)
\psline[linewidth=0.028222222cm,linestyle=dashed,dash=0.16cm 0.16cm](0.84,-1.9317187)(5.84,2.0682812)
\pswedge[linewidth=0.028222222](0.84,-1.9317187){0.6}{0.0}{38.66}
\usefont{T1}{ppl}{m}{n}
\rput(3.315,-2.2317188){$100$ m}
\usefont{T1}{ppl}{m}{n}
\rput(1.9145312,-1.6317188){$38.7^\circ$}
\usefont{T1}{ppl}{m}{n}
\rput(0.6,-1.9317187){$Q$}
\usefont{T1}{ppl}{m}{n}
\rput(5.8145313,-2.2317188){$B$}
\usefont{T1}{ppl}{m}{n}
\rput(5.8145313,2.3682814){$T$}
\usefont{T1}{ppl}{m}{n}
\rput(5.514531,0.06828125){$h$}
\pscircle[linewidth=0.02,dimen=outer](0.15,-1.8619757){0.1}
\psline[linewidth=0.02cm](0.15,-1.9619758)(0.15,-2.3482811)
\psline[linewidth=0.02cm](0.01,-2.5082812)(0.15,-2.3482811)
\psline[linewidth=0.02cm](0.15,-2.3428009)(0.27,-2.5082812)
\psline[linewidth=0.02cm](0.03,-2.1307442)(0.31,-2.1282814)
\psline[linewidth=0.04](5.85,-1.6082813)(5.49,-1.6082813)(5.49,-1.9282813)
\end{pspicture} 
}
\end{center}
\\

% \end{figure}
% \caption{Determining die height of a building using trigonometry.}
% \label{trig:height}

Die diagram gegewe toon 'n gebou van onbekende hoogte $h$. As ons $(Q)$ $100$ m wegstap en die hoek meet van die grond na die toppunt van die gebou $(T)$, vind ons die hoek is $38,{7}^{\circ }$. Hierdie word genoem die elevasiehoek.\\
Ons het 'n reghoekige driehoek en ken die lengte van een sy en 'n hoek. Ons kan dus die hoogte van die gebou bereken (korrek tot die naaste meter.)}
{
\westep{Identifiseer die teenoorstaande-, aanliggende- en die skuinssy}

\westep{}
In $\triangle QTB$,
\begin{equation*}
\begin{array}{ccl}\hfill tan~38,7^{\circ }& =& \frac{\mbox{\footnotesize teenoorstaande sy}}{\mbox{\footnotesize aanliggende sy}}\hfill \vspace{4pt} \\
 & =& \dfrac{h}{100}\hfill
  \end{array}
\end{equation*}
\westep{Herrangskik en los op vir $h$}

\begin{equation*}
\begin{array}{ccl}

\hfill h& =& 100\times tan~38,7^{\circ }\hfill \\
& =& 80\hfill
  \end{array}
\end{equation*}

\westep{Skryf finale antwoord neer}
Die hoogte van die gebou is $80$ m.
}
\end{wex}


\begin{wex}{Elevasie-en depressiehoeke}
{'n Woonstelblok is $200$ m weg van 'n selfoontoring. Iemand staan by $B$. Sy meet die hoek van $B$ na die top van die toring $E$ en vind dit is $34^{\circ}$ (die elevasiehoek). Sy meet dan die hoek vanaf $B$ na die onderpunt van die toring $C$ en vind dit is $62^{\circ}$ (die depressiehoek.)
\\Wat is die hoogte van die selfoontoring (korrek tot die naaste meter)?\\

\\
\begin{center}
\scalebox{1} % Change this value to rescale die drawing.
{
\begin{pspicture}(0,-3.893125)(10.969063,3.893125)
\definecolor{color194}{rgb}{0.2,0.2,0.2}
\definecolor{color353b}{rgb}{0.6,0.6,0.6}
\definecolor{color649b}{rgb}{0.8,0.8,1.0}
\psframe[linewidth=0.04,linecolor=color194,dimen=outer,fillstyle=solid](2.3,0.8496875)(0.18,-3.2103126)
\psframe[linewidth=0.04,linecolor=color194,dimen=outer,fillstyle=solid,fillcolor=color353b](0.62,-0.2903125)(0.4,-0.69031256)
\psframe[linewidth=0.04,linecolor=color194,dimen=outer,fillstyle=solid,fillcolor=color353b](1.32,-0.3103125)(1.1,-0.7103125)
\psframe[linewidth=0.04,linecolor=color194,dimen=outer,fillstyle=solid,fillcolor=color353b](0.62,-1.0303124)(0.4,-1.4303125)
\psframe[linewidth=0.04,linecolor=color194,dimen=outer,fillstyle=solid,fillcolor=color353b](0.62,0.4896875)(0.4,0.0896875)
\psframe[linewidth=0.04,linecolor=color194,dimen=outer,fillstyle=solid,fillcolor=color353b](1.32,0.4896875)(1.1,0.0896875)
\psframe[linewidth=0.04,linecolor=color194,dimen=outer,fillstyle=solid,fillcolor=color353b](2.02,0.4896875)(1.8,0.0896875)
\psframe[linewidth=0.04,linecolor=color194,dimen=outer,fillstyle=solid,fillcolor=color353b](2.02,-0.3103125)(1.8,-0.7103125)
\psframe[linewidth=0.04,linecolor=color194,dimen=outer,fillstyle=solid,fillcolor=color353b](1.32,-1.0303124)(1.1,-1.4303125)
\psframe[linewidth=0.04,linecolor=color194,dimen=outer,fillstyle=solid,fillcolor=color353b](2.02,-1.0303124)(1.8,-1.4303125)
\psframe[linewidth=0.04,linecolor=color194,dimen=outer,fillstyle=solid,fillcolor=color353b](0.62,-1.7303125)(0.4,-2.1303124)
\psframe[linewidth=0.04,linecolor=color194,dimen=outer,fillstyle=solid,fillcolor=color353b](1.32,-1.7303125)(1.1,-2.1303124)
\psframe[linewidth=0.04,linecolor=color194,dimen=outer,fillstyle=solid,fillcolor=color353b](2.02,-1.7303125)(1.8,-2.1303124)
\psframe[linewidth=0.04,linecolor=color194,dimen=outer,fillstyle=solid,fillcolor=color649b](0.5,-2.6103127)(0.2,-3.2103126)
\pscircle[linewidth=0.04,linecolor=color194,dimen=outer,fillstyle=gradient,gradlines=2000,gradmidpoint=1.0](0.35,-2.8603127){0.05}
\psframe[linewidth=0.04,linecolor=color194,dimen=outer,fillstyle=solid,fillcolor=color353b](10.2,2.2896874)(10.0,-3.2103126)
\psline[linewidth=0.04cm,linecolor=color194](10.1,3.2896874)(10.1,2.2896874)
\psline[linewidth=0.025999999cm,linecolor=color194](2.24,0.82968754)(9.98,-3.1703126)
\psline[linewidth=0.024cm,linecolor=color194](2.28,0.8496875)(10.12,3.2896874)
\psdots[dotsize=0.12,linecolor=color194](10.1,3.2896874)
\psline[linewidth=0.024cm,linecolor=color194,linestyle=dashed,dash=0.16cm 0.16cm](2.26,0.82968754)(10.02,0.82968754)
% \usefont{T1}{ptm}{m}{n}
\rput(5.776562,-3.6903126){$200$ m}
\usefont{T1}{ptm}{m}{n}
\rput(2.78,-2.7703125){$ A$}
\usefont{T1}{ptm}{m}{n}
\rput(2.2557812,1.1896875){$B$}
\usefont{T1}{ptm}{m}{n}
\rput(10.574532,-3.1703126){$C$}
\usefont{T1}{ptm}{m}{n}
\rput(10.527812,0.82968754){$D$}
\usefont{T1}{ptm}{m}{n}
\rput(10.371095,3.6896875){$E$}
\usefont{T1}{ptm}{m}{n}
\rput(3.8720312,1.0896875){$34^\circ$}
\usefont{T1}{ptm}{m}{n}
\rput(3.6520312,0.44968748){$62^\circ$}
\psarc[linewidth=0.024,linecolor=color194,arrowsize=0.05291667cm 2.0,arrowlength=1.4,arrowinset=0.4]{->}(4.07,1.0396874){0.49}{336.25052}{58.24052}
\psarc[linewidth=0.024,linecolor=color194,arrowsize=0.05291667cm 2.0,arrowlength=1.4,arrowinset=0.4]{<-}(3.56,0.80968744){0.94}{294.44397}{0.0}
\psframe[linewidth=0.04,linecolor=color194,dimen=outer,fillstyle=solid,fillcolor=color649b](10.38,-3.2103126)(0.0,-3.3903124)
\psline[linewidth=0.04](2.28,-2.886875)(2.6,-2.886875)(2.6,-3.226875)
\psline[linewidth=0.04](9.7,0.51312506)(9.7,1.153125)(10.05669,1.1505005)
\psline[linewidth=0.04cm](9.72,0.53312504)(10.0,0.53312504)
\psline[linewidth=0.04cm](9.72,0.83312505)(10.0,0.83312505)
\end{pspicture} 
}\\
Let op: Die diagram is nie op skaal geteken nie
\end{center}
}
{
\westep{Om die hoogte $CE$ te bepaal, bereken eers die lengtes $DE$ en $CD$}
$\triangle BDE$ en $\triangle BDC$ is beide reghoekige driehoeke. In elk van die driehoeke is die lengte $BD$ bekend, dus kan ons die sye van die driehoeke bereken.

\westep{Bereken $CD$}
Die lengte $AC$ is gegee. $CABD$ is 'n reghoek, dus is $BD = AC = 200\mbox{m}$.
In $\triangle CBD$, 
\begin{eqnarray*}
tan ~C\hat{B}D &=& \frac{CD}{BD}\\
\therefore CD&=&BD\times tan~ C\hat{B}D \\
&=& 200\times tan~ 62^{\circ} \\
&=& 376\mbox{ m}
\end{eqnarray*}

\westep{Bereken $DE$}
In $\triangle DBE$,
\begin{eqnarray*}
tan~ D\hat{B}E &=& \frac{DE}{BD}\\
\therefore DE&=&BD\times tan ~D\hat{B}E \\
&=& 200\times tan~ 34^\circ \\
&=&135\mbox{ m}
\end{eqnarray*}

\westep{Sommeer die twee hoogtes om die finale antwoord te kry} 
Die hoogte van die toring is $CE=CD+DE=135 \mbox{ m}+376\mbox{ m}=511\mbox{ m}$.
}
\end{wex}


\begin{wex}{Gebouplan}
{Mnr Nkosi het 'n motorhuis by sy huis en hy besluit om 'n afdak aan te bou aan die kant van die motorhuis. Die motorhuis is $4$ m hoog en sy plaat vir die dak is $5$ m lank. As die hoek van die dak $5^\circ$ is, hoe hoog moet hy die muur $BD$ bou? Gee die antwoord korrek tot $1$ desimale plek.
\\
\begin{center}
\scalebox{1} % Change this value to rescale the drawing.
{
\begin{pspicture}(0,-1.783125)(9.368906,1.783125)
\definecolor{color194}{rgb}{0.2,0.2,0.2}
\definecolor{color247b}{rgb}{0.8,0.8,0.8}
\psframe[linewidth=0.002,linecolor=white,linestyle=dotted,dotsep=0.16cm,dimen=outer,fillstyle=solid,fillcolor=color247b](8.471094,0.67968744)(8.311094,-1.6403126)
\psframe[linewidth=0.04,linecolor=white,dimen=outer,fillstyle=solid,fillcolor=color247b](4.3910937,0.25968745)(3.8510938,-1.6203126)
\psframe[linewidth=0.04,linecolor=white,dimen=outer,fillstyle=solid,fillcolor=color247b](1.2910937,0.25968745)(0.75109375,-1.6603125)
\psframe[linewidth=0.002,linecolor=white,linestyle=dotted,dotsep=0.16cm,dimen=outer,fillstyle=solid,fillcolor=color247b](4.3910937,1.3196875)(0.79109377,0.19968745)
% \usefont{T1}{ptm}{m}{n}
\rput(2.5167189,-0.20031255){Garage}
% \usefont{T1}{ptm}{m}{n}
\rput(7.78125,1.2396874){Roof}
% \usefont{T1}{ptm}{m}{n}
\rput(8.949843,-0.50031257){Wall}
\psline[linewidth=0.04cm](0.7710937,1.2996874)(0.7710937,-1.6003126)
\psline[linewidth=0.04cm](0.75109375,1.2996874)(4.351094,1.2996874)
\psline[linewidth=0.04cm](4.3710938,1.2996874)(4.3710938,-1.6003126)
\psline[linewidth=0.04cm](0.7710937,-1.6003126)(1.2710937,-1.6003126)
\psline[linewidth=0.04cm](1.2710937,-1.6003126)(1.2710937,0.19968745)
\psline[linewidth=0.04cm](1.2710937,0.19968745)(3.8710938,0.19968745)
\psline[linewidth=0.04cm](3.8710938,0.19968745)(3.8710938,-1.6003126)
\psline[linewidth=0.04cm](3.8710938,-1.6003126)(4.3710938,-1.6003126)
\psline[linewidth=0.024cm,linecolor=color194](4.3710938,1.2996874)(8.44,0.663125)
\psline[linewidth=0.027999999cm,linecolor=color194,linestyle=dashed,dash=0.16cm 0.16cm](8.4,0.663125)(4.2510934,0.6996874)
\psline[linewidth=0.04cm,linecolor=color194](8.291094,0.67968744)(8.291094,-1.6203126)
\psline[linewidth=0.04cm,linecolor=color194](8.451094,0.67968744)(8.451094,-1.6203126)
\psline[linewidth=0.018cm,linecolor=color194](8.46,-1.616875)(0.9710938,-1.6003126)
% \usefont{T1}{ptm}{m}{n}
\rput(0.3,0.07968745){$4$ m}
% \usefont{T1}{ptm}{m}{n}
\rput(6.3395314,1.2596875){$5$ m}
% \usefont{T1}{ptm}{m}{n}
\rput(8.906875,0.6996874){$ B$}
% \usefont{T1}{ptm}{m}{n}
\rput(4.6251564,0.41968745){$A$}
% \usefont{T1}{ptm}{m}{n}
\rput(6.073125,0.8796874){$5^\circ$}
% \usefont{T1}{ptm}{m}{n}
\rput(4.5270314,1.5796875){$C$}
\psdots[dotsize=0.12,linecolor=color194](4.3710938,0.6996874)
\psdots[dotsize=0.12,linecolor=color194](4.3710938,1.2996874)
% \usefont{T1}{ptm}{m}{n}
\rput(8.858906,-1.5803125){$ D$}
\psdots[dotsize=0.1378129,linecolor=color194](8.271093,0.67968744)
\psline[linewidth=0.04](4.38,0.943125)(4.64,0.943125)(4.64,0.703125)(4.64,0.703125)(4.64,0.723125)
\end{pspicture} 
}
\end{center}
}{
\westep{Identifiseer die teenoorstande sy, die aanliggende sy en die skuinssy}
$\triangle ABC$ het 'n regte hoek. Die skuinssy en 'n hoek is bekend, dus kan ons $AC$ bereken. Die hoogte van die muur $BD$ is dan die hoogte van die motorhuis minus $AC$.
\begin{eqnarray*}
 sin~A\hat{B}C &=& \frac{AC}{BC} \\
\therefore AC &=& BC \times sin~A\hat{B}C\\
&=& 5~sin~5^{\circ}\\
&=& 0,4\mbox{ m}\\
\\
\therefore BD&=& 4\mbox{ m}-0,4\mbox{ m}\\
&=& 3,6\mbox{ m}
\end{eqnarray*}


\westep{Skryf die finale antwoord neer}  
Mr Nkosi moet sy muur $3,6$ m hoog bou.
}
\end{wex}

\begin{exercises}{}
{
\begin{enumerate}[noitemsep, label=\textbf{\arabic*}. ] 

\item 'n Seun wat 'n vlieer vlieg staan $30~$m vanaf 'n punt direk onder die vlieer. As die vlieer se tou $50~$m lank is, vind die elevasiehoek van die vlieer.
\item Wat is die elevasiehoek van die son as 'n boom wat $7,15$ m hoog is 'n skaduwee van $10,1$ m het?
\item Susan is $1,4$ m lank. Vanaf 'n afstand van $300$ m kyk sy op na die bopunt van 'n vuurtoring. die elevasiehoek is $5^{\circ}$. Bepaal die hoogte van die vuurtoring tot die naaste meter.
\item 'n Leer met lengte $25$ m rus teen 'n muur teen 'n hoek van $37^{\circ}$. Vind die afstand tussen die muur en die onderpunt van die leer. 

\end{enumerate}

\par \raisebox{-0.2em}{\includegraphics[height=1em]{../icons/www.pdf}} Vind die antwoorde met die kortkodes:
\par \begin{tabular}[h]{cccccc}
(1.) lcY  &  (2.) lcr  & \end{tabular}
    
}
\end{exercises} 



\section{Definisie van verhoudings in die Cartesiese vlak}

Ons het trigonometriese verhoudings gedefinieer deur reghoekige driehoeke te gebruik. Ons kan hierdie definisies uitbrei na enige hoek, deur daarop te let dat die definisies nie afhanklik is van die lengte van die sye van 'n driehoek nie, maar wel net afhanklik van die grootte van die hoek. As ons dus enige punt op 'n Cartesiese vlak plot en 'n lyn trek vanaf die oorsprong na daardie punt kan ons die hoek van die lyn bereken. In die figuur hieronder is punte $P$ and $Q$ reeds geplot. 'n Lyn word getrek vanaf die oorsprong na elke punt. Die stippellyne wys hoe ons reghoekige driehoeke vir elke punt kan konstrueer. Die stippellyn moet altyd na die $x$-as getrek word. Nou kan ons hoeke $A$ and $B$ vind. Ons kan ook die definisies van die omgekeerdes op dieselfde manier uitbrei:



\setcounter{subfigure}{0}
\begin{center}
\scalebox{1} % Change this value to rescale die drawing.
{
\begin{pspicture}(0,-4.0784373)(9.84,4.1184373)
\rput(4.0,-0.0784375){\psaxes[linewidth=0.04,arrowsize=0.05291667cm 2.0,arrowlength=1.4,arrowinset=0.4,ticksize=0.004cm]{<->}(0,0)(-4,-4)(4,4)}

\psline[linewidth=0.04cm,dotsize=0.07055555cm 2.0]{-*}(4.04,-0.0584375)(6.94,1.8615625)
\psline[linewidth=0.04cm,dotsize=0.07055555cm 2.0]{-*}(4.0,-0.0784375)(1.02,1.8815625)
\psline[linewidth=0.04cm,linestyle=dashed,dash=0.16cm 0.16cm](6.96,1.8615625)(6.96,-0.0984375)
\psline[linewidth=0.04cm,linestyle=dashed,dash=0.16cm 0.16cm](1.02,1.8015625)(1.02,-0.0584375)
\psline[linewidth=0.04,fillstyle=solid](9.82,-1.0984375)(9.82,-1.0984375)(9.82,-1.0984375)(9.82,-1.0984375)
\psline[linewidth=0.04,fillstyle=solid](1.02,0.2015625)(1.28,0.2015625)(1.28,-0.0384375)(1.28,-0.0584375)
\psline[linewidth=0.04,fillstyle=solid](6.7,-0.0584375)(6.7,0.1815625)(6.96,0.1815625)(6.96,0.2015625)
\rput{-279.77182}(2.8218253,-2.9484913){\psarc[linewidth=0.04](3.160762,0.20043099){0.48562816}{40.75301}{136.25127}}
\rput{-56.56522}(1.9601717,3.1761396){\psarc[linewidth=0.04](3.9315972,-0.23347138){0.83129716}{66.93879}{193.12923}}
\rput{-72.17863}(3.2203913,4.738851){\psarc[linewidth=0.04](4.860762,0.160431){0.48562816}{44.75812}{136.25127}}
\usefont{T1}{ptm}{m}{n}
\rput(5.0403123,0.1915625){$A$}
\usefont{T1}{ptm}{m}{n}
\rput(4.2,-0.3){$O$}
\usefont{T1}{ptm}{m}{n}
\rput(4.273125,0.7915625){$B$}
\usefont{T1}{ptm}{m}{n}
\rput(2.9954689,0.2115625){$x$}
\usefont{T1}{ptm}{m}{n}
\rput(7.2570314,2.2315626){$P(2;3)$}
\usefont{T1}{ptm}{m}{n}
\rput(0.9845312,2.1915624){$Q(-2;3)$}
\usefont{T1}{ptm}{m}{n}
\rput(8.235469,0.1315625){$x$}
\usefont{T1}{ptm}{m}{n}
\rput(4.275781,3.9915626){$y$}
\end{pspicture} 
} 
\end{center}
Uit die ko-ordinate van $P(2;3)$ weet ons dat die lengte van die sy oorkant $\hat{A}$ $3$ is en die lengte van dei aanliggende sy $2$. Deur $tan~\hat{A}=\frac{\mbox{\footnotesize teenoorstaande}}{\mbox{\footnotesize aanliggend}} = \frac{3}{2}$ te gebruik, bvereken ons dat $\hat{A}=56,3^{\circ}$ is.\par

Ons kan ook die Teorema van Pythagoras gebruik om die skuinssy van die driehoek te bereken, dan $\hat{A}$ bereken deur $sin~\hat{A} = \frac{\mbox{\footnotesize teenoorstaande}}{\mbox{\footnotesize skuins}}$ or $cos~\hat{A} = \frac{\mbox{\footnotesize aanliggend}}{\mbox{\footnotesize skuins}}$\par
Beskou punt $Q(-2;3)$. Ons definieer $\hat{B}$ as die hoek tussen die lyn $OQ$ en die positiewe $x$-axis. Hierdie word die standaard hoekposisie genoem. Laat $\hat{x}$ die hoek wees tussen lyn $OQ$ en die negatiewe $X$-as sodat $\hat{B} + \hat{x} = 180^{\circ}$.
\par

Uit die ko-ordinate $Q(-2;3)$ weet ons die lengte van die sy oorkant $\hat{x}$ is $3$ en die lengte van die aanliggende sy is $2$. Deur $tan~\hat{x}=\frac{\mbox{\footnotesize teenoorstaande}}{\mbox{\footnotesize aanliggend}} = \frac{3}{2}$ bereken ons $\hat{x}=56,3^{\circ}$.
\\Dus $\hat{B}=180^{\circ} - \hat{x} = 123,7^{\circ}$.\par

Soortgelyk, 'n alternatiewe metode is om die skuinssy te bereken deur die Teorema van Pythagoras ete bereken, dan $\hat{x}$ using $sin~\hat{x} = \frac{\mbox{\footnotesize teenoorstaande}}{\mbox{\footnotesize skuins}}$ or $cos~\hat{x} = \frac{\mbox{\footnotesize aanliggend}}{\mbox{\footnotesize skuins}}$ te gebruik. \par
Sou ons 'n sirkel trek, gesentreerd om die oorsprong $(0)$ en deur punt $P$, dan is die lengte van die oorsprong tot by punt $P$ die radius van die sirkel, wat ons aandui met $r$. Ons kan al die trigonometriese funksies dan herskryf in terme van $x$, $y$ and $r$.

Die algemene definisies van die trigonometriese funksies is:

\begin{equation*}
\begin{array}{cccccc}\hfill sin~\theta & =& \frac{y}{r}\hfill & cosec~\theta & =& \frac{r}{y}\hfill \\
 \hfill cos~\theta & =& \frac{x}{r}\hfill & sec~\theta & =& \frac{r}{x}\hfill \\
 \hfill tan~\theta & =& \frac{y}{x}\hfill & cot~\theta & =& \frac{x}{y}\hfill \end{array}
\end{equation*}
\Note{In trigonometrie word hoeke altyd gemeet vanaf die positiewe $x$-axis in 'n antiklokse rigting.}

Die karetesies vlak is verdeel in $4$ kwadrante in 'n anti-klokse rigting soos aangedui in die diagram hieronde. Let daarop dat $r$ altyd positief i, maar dat die waardes van $x$ en $y$ verander afhangende van die posisie van die punt op die Cartesiese vlak. As 'n resultaat kan die trigonometriese 
verhoudings posiief of negatief wees. Die letters $C$, $A$, $S$ en $T$ dui aan watter van die funksies is positief in elke kwadrant:\\
\\
Kwadrant I: Alle verhoudings is positief.\\
Kwadrant II: $y$ waardes is positief, dus is $sin$ and $cosec$ positief.\\
Kwadrant III: Beide $x$ en $y$ waardes is negatief, dus is $tan$ en $cot$ positief. \\
Kwadrant IV: $x$ waardes is positief, dus is $cos$ en $sec$ positief.\par


\Note{$r$ is 'n lengte en dus altyd positief.}
\begin{center}
\scalebox{1} % Change this value to rescale die drawing.
{
\begin{pspicture}(0,-4.1584377)(8.42125,4.1984377)
\rput(4.0,-0.1584375){\psaxes[linewidth=0.04,arrowsize=0.05291667cm 2.0,arrowlength=1.4,arrowinset=0.4,labels=none,ticks=none,ticksize=0.10583333cm]{<->}(0,0)(-4,-4)(4,4)}
% \usefont{T1}{ptm}{m}{n}
\rput(6.0853124,2.2715626){Quadrant I}
% \usefont{T1}{ptm}{m}{n}
\rput(1.7453125,2.3115625){Quadrant II}
% \usefont{T1}{ptm}{m}{n}
\rput(1.8053125,-2.3084376){Quadrant III}
% \usefont{T1}{ptm}{m}{n}
\rput(6.3515625,-2.2684374){Quadrant IV}
% \usefont{T1}{ptm}{m}{n}
\rput(4.7403126,0.8315625){A }
% \usefont{T1}{ptm}{m}{n}
\rput(4.7723436,0.4315625){all}
% \usefont{T1}{ptm}{m}{n}
\rput(3.2923439,0.8515625){S}
% \usefont{T1}{ptm}{m}{n}
\rput(3.2390625,0.4715625){$sin$}
% \usefont{T1}{ptm}{m}{n}
\rput(3.2109375,-0.8484375){T}
% \usefont{T1}{ptm}{m}{n}
\rput(3.1639063,-1.1684375){$tan$}
% \usefont{T1}{ptm}{m}{n}
\rput(4.7857814,-0.8484375){C}
% \usefont{T1}{ptm}{m}{n}
\rput(4.815,-1.1484375){$cos$}
\usefont{T1}{ptm}{m}{n}
\rput(8.275469,-0.1084375){$x$}
\usefont{T1}{ptm}{m}{n}
\rput(4.1557813,4.0715623){$y$}
\usefont{T1}{ptm}{m}{n}
\rput(7.810625,0.0715625){$0^{\circ}$}
\usefont{T1}{ptm}{m}{n}
\rput(7.60875,-0.4484375){$360^{\circ}$}
\usefont{T1}{ptm}{m}{n}
\rput(4.3410935,3.5515625){$90^{\circ}$}
\usefont{T1}{ptm}{m}{n}
\rput(0.32265624,0.0715625){$180^{\circ}$}
\usefont{T1}{ptm}{m}{n}
\rput(3.529375,-3.9084375){$270^{\circ}$}

\usefont{T1}{ptm}{m}{n}
\rput(4.170625,-0.4084375){$0$}
\end{pspicture} 
}
\end{center}
\\
 Hierdie diagram staan bekend as die CAST diagram.
\begin{wex}{Verhoudings in die Cartesiese vlak}
{$P(-3;4)$ is a punt in die Cartesiese vlak. $X\hat{O}P=\theta$. Sonder om 'n sakrekenaar te gebruik, bepaal die waarde van: 
\begin{enumerate}[noitemsep, label=\textbf{\arabic*}. ] 
\item $cos~\theta$
\item $3~tan~\theta$
\item $\frac{1}{2}~cosec~\theta$
\end{enumerate}
}
{
\westep{Skets punt point $P$ in die Cartesiese vlak en merk dit $\theta$} 
\begin{center}
\scalebox{1} % Change this value to rescale die drawing.
{
\begin{pspicture}(0,-3.1584375)(6.44125,3.1984375)
\rput(3.0,-0.1584375){\psaxes[linewidth=0.04,arrowsize=0.05291667cm 2.0,arrowlength=1.4,arrowinset=0.4,labels=none,ticks=none,ticksize=0.10583333cm]{<->}(0,0)(-3,-3)(3,3)}
\usefont{T1}{ptm}{m}{n}
\rput(6.295469,0.0115625){$x$}
\usefont{T1}{ptm}{m}{n}
\rput(3.1757812,3.0715625){$y$}
\psline[linewidth=0.04cm,dotsize=0.07055555cm 2.0]{-*}(3.02,-0.1784375)(1.6,1.8615625)
\usefont{T1}{ptm}{m}{n}
\rput(1.3670312,2.2715626){$P(-3;4)$}
\rput{-45.048485}(0.95434314,2.1533978){\psarc[linewidth=0.04](3.073445,-0.073917985){0.6961048}{36.728504}{180.0}}
\usefont{T1}{ptm}{m}{n}
\rput(3.150625,-0.4484375){$0$}
\usefont{T1}{ptm}{m}{n}
\rput(3.2709374,0.1915625){$\theta$}
\end{pspicture} 
}
\end{center}

\westep{Gebruik Pythagoras om $r$ te bereken}
\begin{equation*}
 \begin{array}{ccl}
    \hfill r^{2} &= & x^{2} + y^{2} \hfill \\
\hfill  &=& (-3)^{2} + (4)^{2} \hfill \\
\hfill  &=& 25 \hfill \\
\hfill r &=& \pm 5\hfill \\
\hfill \therefore r &=& 5 \hfill \\ 
\end{array}
\end{equation*}
\westep{Vervang waardes vir $x$, $y$ en $r$ in die nodige verhoudings}
\begin{enumerate}[itemsep=5pt, label=\textbf{\arabic*}. ] 
   \item $cos~\theta = \dfrac{x}{r} = -\dfrac{3}{5}$
\item $3~tan~\theta = 3\Big(\dfrac{y}{x}\Big) = 3\Big(\dfrac{4}{-3}\Big) = -4 $
\item $\dfrac{1}{2}~cosec~\theta = \dfrac{1}{2}\Big(\dfrac{r}{y}\Big) = \dfrac{1}{2}\Big(\dfrac{5}{4}\Big) = \dfrac{5}{8} $
  \end{enumerate}
}
\end{wex}


\begin{wex}{Verhoudings in die Cartesiese vlak}{
$X\hat{O}K = \theta$ is 'n hoek in die derde kwadrant en $K$ is die punt $(-5;y)$. $OK$ is $13$ eenhede. Bepaal sonder 'n rekenaar:
  \begin{enumerate}[noitemsep, label=\textbf{\arabic*}. ] 
   \item Die waarde van $y$
\item Bewys dat $tan^{2}~\theta + 1 = sec^{2}~\theta$
  \end{enumerate}
}
{
\westep{Teken punt $K$ in die Cartesiese vlak merk dit $\theta$} 
\begin{center}
\scalebox{1} % Change this value to rescale die drawing.
{
\begin{pspicture}(0,-3.1584375)(6.44125,3.1984375)
\rput(3.0,-0.1584375){\psaxes[linewidth=0.04,arrowsize=0.05291667cm 2.0,arrowlength=1.4,arrowinset=0.4,labels=none,ticks=none,ticksize=0.10583333cm]{<->}(0,0)(-3,-3)(3,3)}
\usefont{T1}{ptm}{m}{n}
\rput(6.295469,0.0115625){$x$}
\usefont{T1}{ptm}{m}{n}
\rput(3.1757812,3.0715625){$y$}
\psline[linewidth=0.04cm,dotsize=0.07055555cm 2.0]{-*}(2.963227,-0.18549357)(1.6,-2.1384375)
\usefont{T1}{ptm}{m}{n}
\rput(1.2170312,-2.5084374){$K(-5;y)$}
\rput{-45.048485}(1.0897402,2.0403726){\psarc[linewidth=0.04](3.0048735,-0.2936738){0.7345531}{56.389397}{272.69913}}
\usefont{T1}{ptm}{m}{n}
\rput(3.150625,-0.4484375){$0$}
\usefont{T1}{ptm}{m}{n}
\rput(2.4109375,0.5115625){$\theta$}
\usefont{T1}{ptm}{m}{n}
\rput(1.9303125,-1.1484375){$13$}
\end{pspicture} 
}
\end{center}
\westep{Use Pythagoras to calculate $y$}
\begin{equation*}
 \begin{array}{ccl}
    \hfill r^{2} &= & x^{2} + y^{2} \hfill \\
\hfill y^{2} &=& r^{2} - x^{2}\hfill \\
\hfill  &=& (13)^{2} - (-5)^{2} \hfill \\
\hfill  &=& 169-25 \hfill \\
\hfill  &=& 144 \hfill \\
\hfill  y&=& \pm 12\hfill 

\end{array}
\end{equation*}
Gegee dat $\theta$ in die derde kwadrant val, dan moet $y$ negatief wees.\\
\begin{equation*}
 \therefore y = -12
\end{equation*}

\westep{Vervang waardes vir $x$, $y$ en $r$ en vereenvoudig}
\begin{equation*}
\begin{array}{rclrcl}
 \hfill \mbox{LHS} &=&  tan^{2}~\theta + 1 	&	\mbox{RHS}	&=&  	sec^{2}~\theta \hfill \\
\hfill &=& \Big(\dfrac{y}{x}\Big)^{2} + 1  	&	&=&  Big(\dfrac{r}{x}\Big)^{2} \hfill \vspace{5pt} \\
\hfill&=&  \Big(\dfrac{-12}{-5}\Big)^{2} + 1 	&	&=&  \Big(\dfrac{13}{-5}\Big)^{2} \hfill \vspace{5pt}\\
\hfill&=&  \Big (\dfrac{144}{25}\Big) + 1  	&	&=&  \Big(\dfrac{169}{25}\Big)^{2} \hfill \vspace{5pt}\\
\hfill &=& \dfrac{144 + 25}{25} 		&	&=&  \dfrac{169}{25} \hfill \vspace{5pt}\\
\hfill &=& \dfrac{169}{25}  			&	&=&  \dfrac{169}{25} \hfill \vspace{5pt}\\
\hfill \therefore \mbox{LHS} &=& \mbox{RHS}&& \hfill

\end{array}
\end{equation*}

}
\end{wex}

\Tip{Maak altyd 'n skets wanneer jy trigonometriese probleme sonder 'n rekenaar moet oplos.}

\begin{exercises}{}
{
  \begin{enumerate}[itemsep=5pt, label=\textbf{\arabic*}. ]
   \item $B$ is a punt in die Cartesiese vlak. Bepaal sonder 'n rekenaar:
% \begin{center}
% \scalebox{1} % Change this value to rescale die drawing.
% {
% \begin{pspicture}(0,-3.1584375)(6.44125,3.1984375)
% \rput(3.0,-0.1584375){\psaxes[linewidth=0.04,arrowsize=0.05291667cm 2.0,arrowlength=1.4,arrowinset=0.4,labels=none,ticks=none,ticksize=0.10583333cm]{<->}(0,0)(-3,-3)(3,3)}
% \usefont{T1}{ppl}{m}{n}
% \rput(6.295469,0.0115625){$x$}
% \usefont{T1}{ppl}{m}{n}
% \rput(3.1757812,3.0715625){$y$}
% \rput(3.4, -0.4){$O$}
% \psline[linewidth=0.04cm,fillcolor=black,dotsize=0.07055555cm 2.0]{-*}(3.003227,-0.16549358)(3.78,-1.8384376)
% \usefont{T1}{ppl}{m}{n}
% \rput(4.577031,-1.9084375){$B(1;-3)$}
% \rput{-45.048485}(1.0590175,2.0862904){\psarc[linewidth=0.04](3.0448735,-0.2336738){0.7345531}{50.714928}{340.5804}}
% \usefont{T1}{ppl}{m}{n}
% \rput(2.393125,0.5115625){$\beta$}
% \end{pspicture} 
% }
% \end{center}
\begin{enumerate}[noitemsep, label=\textbf{(\alph*)} ]
 \item $OB$
\item $cos~\beta$
\item $cosec~\beta$
\item $tan~\beta$
\end{enumerate}


\item As $sin~\theta= 0,4$ en $\theta$ is 'n stomp hoek, bepaal:
% \begin{center}
% \scalebox{1} % Change this value to rescale die drawing.
% {
% \begin{pspicture}(0,-3.1584375)(6.44125,3.1984375)
% \rput(3.0,-0.1584375){\psaxes[linewidth=0.04,arrowsize=0.05291667cm 2.0,arrowlength=1.4,arrowinset=0.4,labels=none,ticks=none,ticksize=0.10583333cm]{<->}(0,0)(-3,-3)(3,3)}
% \usefont{T1}{ppl}{m}{n}
% \rput(6.295469,0.0115625){$x$}
% \usefont{T1}{ppl}{m}{n}
% \rput(3.1757812,3.0715625){$y$}
% \psline[linewidth=0.04cm,fillcolor=black](3.0,-0.1784375)(0.8,0.7415625)
% \usefont{T1}{ppl}{m}{n}
% \rput(0.7292187,-0.4884375){$(x;2)$}
% \rput{-45.048485}(1.0440966,2.0498228){\psarc[linewidth=0.04](2.9934452,-0.23391798){0.6961048}{50.894093}{199.87982}}
% \usefont{T1}{ppl}{m}{n}
% \rput(3.150625,-0.4484375){$0$}
% \usefont{T1}{ppl}{m}{n}
% \rput(3.6509376,0.5915625){$\theta$}
% \psline[linewidth=0.04cm,linestyle=dashed,dash=0.16cm 0.16cm](0.8,0.7215625)(0.8,-0.1784375)
% \usefont{T1}{ptm}{m}{n}
% \rput(0.5234375,0.2315625){$2$}
% \end{pspicture} 
% }
% \end{center}
\begin{enumerate}[noitemsep, label=\textbf{(\alph*)} ]
% \setcounter{enumi}{4}
 \item $cos~\theta$
\item $\sqrt{21}~tan~\theta$
\end{enumerate}
\end{enumerate}
}
\end{exercises}

\summary
\begin{itemize}[noitemsep]
\item Ons kan drie trigonometriese verhoudings definieer vir reghoekige driehoeke: sinus ($sin$), cosinus ($cos$) en tangens ($tan$).
\item Elk van hierdie verhoudings het 'n omgekeerde: cosecant ($cosec$), secant ($sec$) and cotangens ($cot$).
\item Ons kan die beginsels van oplossing van vergelykings en die trigonometriese verhoudings gebruik om eenvoudige trigonometriese probleme op te los.
\item Ons kan probleme in twee dimensies oplos met reghoekige driehoeke.
\item Vir sekere spesiale hoeke kan ons die waardes van $sin$, $cos$ and $tan$ maklik bepaal sonder 'n rekenaar.
\item Ons kan die definisies van die trigonometriese funksies uitbrei na enige hoek.
\item Trigonometrie word gebruik om probleme in twee dimensies op te los, soos om die hoogte van 'n gebou te bepaal.
\end{itemize}


\begin{eocexercises}{}

\begin{enumerate}[itemsep=6pt, label=\textbf{\arabic*}. ] 
\item Sonder om 'n sakrekenaar te gebruik, bepaal die waarde van
\begin{equation*}
sin~60^{\circ}~cos~30^{\circ}-cos~60^{\circ}sin~30^{\circ} + tan~45^{\circ}
\end{equation*}
\item As $3~tan~\alpha = -5$ en $0^{\circ} < \alpha < 270^{\circ}$, gebruik 'n skets om te bepaal:
    \begin{enumerate}[noitemsep, label=\textbf{(\alph*)} ]
    \item $cos~\alpha$
    \item $tan^{2}~\alpha - sec^{2}~\alpha$
    \end{enumerate}
\item Los op vir $\theta$ as $\theta$ 'n positiewe skerp hoek is:
    \begin{enumerate}[noitemsep, label=\textbf{(\alph*)} ]
    \item $2~sin~\theta = 1,34$
    \item $1 - tan~\theta = -1$
    \item $cos~2\theta = sin~40^{\circ}$ 
    \item $\dfrac{sin~\theta}{cos~\theta}= 1$
    \end{enumerate}


\item Bereken die onbekende lengtes in die diagramme hieronder:
\begin{center}
\scalebox{1}  
{ 
\begin{pspicture}(0,-2.0390613)(8.035,2.0253136) 
\psline[linewidth=0.04cm](0.02,0.02093862)(3.06,0.02093862) 
\psline[linewidth=0.04cm](0.02,0.02093862)(0.02,-2.0190613) 
\psline[linewidth=0.04cm](0.04,-1.9990613)(3.04,0.02093862) 
\psframe[linewidth=0.04,dimen=outer](0.26,0.04093862)(0.0,-0.21906137) 
\psline[linewidth=0.04cm](0.9944844,1.4531701)(3.04,0.04093862) 
\psline[linewidth=0.04cm](0.96,1.4209386)(0.024381146,0.02542873) 
\rput{-33.90198}(-0.5512336,0.79249614){\psframe[linewidth=0.04,dimen=outer](1.1544285,1.4305158)(0.8944285,1.1705158)} 
\psline[linewidth=0.04cm](2.2384837,1.992771)(3.04,0.04093862) 
\psline[linewidth=0.04cm](2.238,1.9929386)(1.0103352,1.4765087) 
\rput{-66.64198}(-0.3803531,3.117777){\psframe[linewidth=0.04,dimen=outer](2.3111107,1.9781737)(2.0511107,1.7181737)} 
% \usefont{T1}{ptm}{m}{n} 
\rput(2.2601562,-0.21406138){$30^{\circ}$} 
% \usefont{T1}{ptm}{m}{n} 
\rput(2.3034375,0.22593862){$25^{\circ}$} 
% \usefont{T1}{ptm}{m}{n} 
\rput(2.4634376,0.77){$20^{\circ}$} 
% \usefont{T1}{ptm}{m}{n} 
\rput(1.8576562,-1.2540613){$16$ cm} 
% \usefont{T1}{ptm}{m}{it} 
\rput(1.3398438,0.20593862){$a$} 
% \usefont{T1}{ptm}{m}{it} 
\rput(1.9,1.1){$b$} 
% \usefont{T1}{ptm}{m}{it} 
\rput(1.58,1.9059386){$c$} 
\psline[linewidth=0.04cm](6.689009,-1.6241995)(6.0199666,-1.4501117) 
\rput{-90.46057}(4.6635084,7.686732){\psframe[linewidth=0.04,dimen=outer](6.2543488,1.6402806)(6.0343485,1.4202806)} 
\rput{-90.46057}(4.4806204,7.508202){\psframe[linewidth=0.04,dimen=outer](6.0743546,1.6417276)(5.8543544,1.4217275)} 
\rput{-104.9971}(9.7950115,4.4890895){\psframe[linewidth=0.04,dimen=outer](6.7298956,-1.4036404)(6.509896,-1.6236403)} 
\psline[linewidth=0.04cm](7.375193,1.6303898)(4.2549725,1.6154709) 
\psline[linewidth=0.04cm](6.055236,1.6410005)(6.030317,-1.4588993) 
\psline[linewidth=0.04cm](6.0304775,-1.4389)(4.2748113,1.5953108) 
\psline[linewidth=0.04cm](7.375193,1.6303898)(6.030317,-1.4588993) 
\psline[linewidth=0.04cm](7.375193,1.6303898)(6.689009,-1.6241995) 
% \usefont{T1}{ptm}{m}{it} 
\rput(4.880625,0.005938619){$d$} 
% \usefont{T1}{ptm}{m}{it} 
\rput(6.7414064,1.8259386){$e$} 
% \usefont{T1}{ptm}{m}{n} 
\rput(5.2226562,1.8259386){$5$ m} 
% \usefont{T1}{ptm}{m}{it} 
\rput(6.2709374,-1.7940614){$f$} 
% \usefont{T1}{ptm}{m}{it} 
\rput(7.2148438,-0.07406138){$g$} 
% \usefont{T1}{ptm}{m}{n} 
\rput(4.756875,1.3859386){ $50^{\circ}$} 
% \usefont{T1}{ptm}{m}{n} 
\rput(6.9460936,1.4059386){$60^{\circ}$} 
% \usefont{T1}{ptm}{m}{n} 
\rput(6.45,-1.2740613){$80^{\circ}$} 
\end{pspicture} 
}
\end{center}

\item In $\triangle PQR$, $PR=20$ cm, $QR=22$ cm and $P\hat{R}Q = 30^{\circ}$. Die loodlyn vanaf $P$ na $QR$ kruis $QR$ by $X$. Bereken 
\begin{enumerate}[noitemsep, label=\textbf{(\alph*)} ]
\item die lengte $XR$ 
\item die lengte $PX$
\item die hoek $Q\hat{P}X$ 
\end{enumerate} 
\item 'n leer met lengte $15$ m rus teen 'n muur. Die voet van die leer is $5$ m from die wall. Vind die hoek tussen die muur en die leer.
\item In die volgende driehoek, vind die hoek $A\hat{B}C$:
\begin{center}
\begin{pspicture}(0,-2.4701562)(5.49875,2.4701562) 
\pspolygon[linewidth=0.04](0.1665625,-1.7301563)(3.3665626,1.9698437)(5.1665626,-1.7301563)(4.1665626,-1.7301563) 
\psline[linewidth=0.04cm](3.3665626,1.9698437)(3.3665626,-1.7301563) 
\rput(3.3871875,2.2798438){$A$} 
\rput(5.3459377,-2.0201561){$B$} 
\rput(3.371875,-2.0201561){$C$} 
\rput(0.07546875,-2.0201561){$D$} 
\rput(3.6,0){$9$} 
\rput(2.7525,-2.3201563){$17$} 
\psline[linewidth=0.04cm,arrowsize=0.05291667cm 2.0,arrowlength=1.4,arrowinset=0.4]{->}(3.0665624,-2.3301563)(5.2665625,-2.3301563) 
\psline[linewidth=0.04cm,arrowsize=0.05291667cm 2.0,arrowlength=1.4,arrowinset=0.4]{->}(2.4665625,-2.3301563)(0.0665625,-2.3301563) 
\psline[linewidth=0.04cm](3.3665626,-1.5301563)(3.5665624,-1.5301563) 
\psline[linewidth=0.04cm](3.5665624,-1.5301563)(3.5665624,-1.7301563) 
\rput(0.8,-1.48){$41^{\circ}$} 
\end{pspicture} 
\end{center}
\item In die volgende driehoek, vind die lengte van sy $CD$:
\begin{center}
\begin{pspicture}(0,-2.2234375)(6.091875,2.2234375) 
\pspolygon[linewidth=0.04](0.1665625,-1.776875)(5.1665626,-1.776875)(5.1665626,1.823125) 
\psline[linewidth=0.04cm](3.4665625,-1.776875)(5.1665626,1.823125) 
\rput(5.2871876,2.033125){$A$} 
\rput(5.3459377,-2.066875){$B$} 
\rput(3.471875,-2.066875){$C$} 
\rput(0.07546875,-2.066875){$D$} 
\rput(5.4,0){$9$} 
\rput(4.490156,0.95){$15^{\circ}$} 
\rput(3.960156,-1.5){$35^{\circ}$} 
\psline[linewidth=0.04cm](4.9665626,-1.576875)(5.1665626,-1.576875) 
\psline[linewidth=0.04cm](4.9665626,-1.576875)(4.9665626,-1.776875) 
\end{pspicture}
\end{center} 
\item $A(5;0)$ en $B(11;4)$. Vind die hoek tussen die lyn deur $A$ en $B$ en die $x$-as. 
\item $C(0;-13)$ en $D(-12;14)$. Vind die hoek tussen die lyn deur $C$ en $D$ en die $y$-as. 


\item 'n Reghoekige driehoek het 'n skuinssy van $13$ mm. Vind die lengte van die ander twee sye as een van die hoeke van die driehoek $50^{\circ}$ is.
\item Een van die hoeke van 'n rombus met omtrek $20$ cm is $30^{\circ}$. 
\begin{enumerate}[noitemsep, label=\textbf{(\alph*)} ]
\item Vind die kante van die rombus. 
\item Vind die lengte van beide diagonale. 
\end{enumerate} 
\item Kaptein Jack seil na 'n krans met 'n hoogte van $10$ m. 
\begin{enumerate}[noitemsep, label=\textbf{(\alph*)} ] 
\item As die afstand vanaf die boot na die bokant van die krans $30$ m is, bereken die elevaasiehoek vanaf die boot na die bokant van die krans (korrek tot tot die naaste heelgetal).
\item As die boat $7$ m nader aan die krans seil, wat is die nuwe elevasiehoek vanaf die boot na die bokant van die krans? 
\end{enumerate} 
\item Gegee die punte $E(5;0)$; $F(6;2)$ en $G(8;-2)$. Vind die hoek $F\hat{E}G$. 
\item  'n Driehoek met hoeke $40^{\circ}$; $40^{\circ}$ en $100^{\circ}$ het 'n omtrek van $20$ cm. Vind die lengte van elke sy  van die driehoek. 

\end{enumerate}
\end{eocexercises}




















